\documentclass[12pt, a4paper, oneside]{report}
\input{preamble}

\begin{document}

% ==========================================================
% TITLE PAGE
% ==========================================================
\begin{titlepage}
\centering
\vspace*{10mm}
{\large INNOPOLIS UNIVERSITY}\\[2mm]
{\normalsize Bachelor of Science in Computer Science}\\[30mm]

{\LARGE\bfseries ViKi: Video-to-Kinematics\\[3mm]
for Object-Centric Retargeting of Human\\[3mm]
Manipulation Demonstrations to Robot Manipulators\par}
\vspace{15mm}

{\large Bachelor's Graduation Thesis}\\[25mm]

\begin{tabular}{rl}
\textit{Author:} & Artem Tuzov\\[4mm]
\textit{Supervisor:} & Alexei Kornaev\\[4mm]
\textit{Group:} & BS23-RO-01\\
\end{tabular}

\vfill
{\normalsize Innopolis, Russia\\ \the\year}
\end{titlepage}

\pagenumbering{roman}

% ==========================================================
% ABSTRACT (EN)
% ==========================================================
\chapter*{Abstract}
\addcontentsline{toc}{chapter}{Abstract}

ViKi is an offline research pipeline that transforms multi-view RGB-D
recordings of tabletop manipulation into hand and object tracks and
candidate robot trajectories. Its intended local task is a blue-cube
transfer between two support cubes using two Azure Kinect DK cameras;
the target embodiment is a simulated old-generation UR10 with a
two-finger gripper. A geometric robot plan is not itself evidence of
a physically successful manipulation.

The reproducible quantitative evaluation in this thesis uses a separate
eight-view HO-Cap pick-and-place cohort: 28 sequences and 29,706 frames.
Prepared hand tracks achieve 10.79~mm matched-joint-weighted MPJPE over
623,826 annotated joint slots. For 112 named object tracks, the macro-mean
relative translation and rotation drift are 5.44~mm and 7.58 degrees after
a first-frame reference alignment, conditional on published instance
masks. Severe orientation outliers remain. The reference annotations
are semi-automatic, and these scores do not validate the local Kinect
rig or absolute object localization. Earlier local cube recordings
and the artifacts derived from them are unavailable; their purported
capture, retargeting, simulation and learning results are excluded.
No physical-robot success, policy improvement or Real-to-Sim-to-Real
transfer is claimed.

\vspace{3mm}
\noindent\textbf{Keywords:} imitation learning; motion retargeting; inverse
kinematics; RGB-D calibration; object-centric representation; robot learning
from human video.

% ==========================================================
% ABSTRACT (RU)
% ==========================================================
\chapter*{Аннотация}
\addcontentsline{toc}{chapter}{Аннотация}
\begin{russian}

ViKi --- исследовательский офлайн-пайплайн преобразования
многокамерных RGB-D записей настольной манипуляции в траектории кисти
и объектов и кандидаты траекторий робота. Планируемая локальная задача
--- перенос синего кубика между двумя опорными кубиками с записью
двумя Azure Kinect DK; целевой робот --- моделируемый UR10 с
двухпалым захватом. Геометрическая траектория сама по себе не
доказывает успешную манипуляцию в физической симуляции.

Воспроизводимая численная оценка в работе выполнена на отдельном
наборе HO-Cap: 28 последовательностей, 29\,706 кадров и восемь камер.
Взвешенная ошибка суставов кисти относительно полуавтоматической
разметки составила 10,79~мм при 623\,826 сопоставленных суставах.
Для 112 поименованных треков объектов средняя по трекам ошибка
относительного смещения составила 5,44~мм, поворота --- 7,58 градуса
после совмещения по первому кадру при использовании опубликованных
масок. Сильные выбросы ориентации сохранены. Эти оценки не проверяют
точность локального рига Kinect или абсолютную локализацию объектов.
Прежние локальные записи с кубиками и производные артефакты
недоступны для проверки; связанные с ними результаты записи,
ретаргетинга, симуляции и обучения исключены. Работа не заявляет
успеха на физическом роботе, улучшения политики или переноса
Real-to-Sim-to-Real.

\vspace{3mm}
\noindent\textbf{Ключевые слова:} обучение с имитацией; ретаргетинг движения;
обратная кинематика; калибровка RGB-D; объектно-центричное представление;
обучение роботов по видео человека.
\end{russian}

% ==========================================================
\chapter*{Acknowledgements}
\addcontentsline{toc}{chapter}{Acknowledgements}
The author thanks Alexei Kornaev for supervision and
guidance throughout this work, and the Innopolis Robotics Society for
maintaining a parallel development fork of the system, whose independent
implementation attempts informed several of the architectural decisions
reported here.

\tableofcontents
\listoffigures
\listoftables

\cleardoublepage
\pagenumbering{arabic}
\setcounter{page}{1}

\chapter{Introduction}
\label{ch:intro}

\section{Background}
\label{sec:background}

Learning manipulation from demonstrations is limited by the cost of acquiring
robot-executable examples. Teleoperation records robot states and actions
directly, but needs a robot and a control interface during collection. Passive
RGB-D capture instead records a person's unconstrained motion with cameras.
It is cheaper to collect, but it does not itself provide robot actions: hand
motion must be reconstructed, related to objects, and transferred to a robot
with different geometry. Even a plausible joint trajectory need not lift or
place an object when contacts and dynamics are simulated.

This thesis studies that chain as a source of training data. The target
experiment is deliberately narrow: a person transfers a blue $40$~mm cube
from a red support cube to a green target cube; a simulated old-generation
UR10 with a Robotiq 2F-85 gripper attempts the same task. The goal is not to
claim a general-purpose manipulation policy, nor to claim deployment on a
physical robot that is not available for this study.

\section{Problem Statement}
\label{sec:problem}

The central question is whether passive demonstrations provide a measurable
advantage when learning a robot policy in physics simulation. Three distinct
gaps intervene. Perception must recover consistent hand and object motion from
two imperfectly calibrated RGB-D views. Retargeting must turn a human motion
into a reachable robot plan with a meaningful gripper command. Finally, the
plan must survive physical execution: a kinematically valid approach may still
miss a cube, slip during transfer, or fail on release.

Object-relative trajectories offer a way to transfer demonstrations to new
cube placements, but they require a reliable object pose and must be checked
again after transfer. Synthetic variations are therefore candidates for
training, not valid demonstrations merely because a geometric transform can
be computed. A second distinction is between learning with object coordinates
and learning from images. The former isolates the quality of transferred data
and control; the latter also tests perception and the visual simulation gap.

\section{Research Gap}
\label{sec:gap}

Prior work covers parts of this problem. Teleoperation research develops
retargeting objectives, while passive-video systems estimate human and object
motion. MimicGen demonstrates that object-centric subtask transfer plus
simulation filtering can expand a small source dataset \cite{mimicgen}. This
work does not claim those operators as novel. Its question is whether an
accessible, camera-only capture pipeline can produce useful robot-learning
examples for the specified cube-transfer task, and which stage limits the
yield. The comparison must include a controller that uses no ViKi examples;
otherwise success in simulation cannot be attributed to the capture pipeline.

\section{Research Questions and Hypotheses}
\label{sec:rq}

\noindent\textbf{RQ.} \emph{Can passive two-view RGB-D demonstrations processed by
\viki\ improve learning of a simulated UR10 cube-transfer policy, compared
with the same training task without ViKi-derived examples?}

\begin{enumerate}[label=\textbf{RQ\arabic*.}, leftmargin=2.2em]
  \item What fraction of recorded episodes survives hand reconstruction,
        object tracking, retargeting and physical replay, and where are the
        dominant losses?
  \item Does object-aware, physics-validated transfer to new cube placements
        improve held-out task success over equal-size world-frame or hand-only
        training sets?
  \item Under matched observations, control, training budget and evaluation
        placements, do ViKi-derived examples improve policy success or sample
        efficiency relative to learning without those examples?
  \item Can a visual policy using simulated camera frames and robot state
        retain task success without privileged object coordinates at test time?
\end{enumerate}

These are testable hypotheses, not findings. H1 predicts a nonzero yield of
physically successful transferred episodes. H2 predicts that object-aware
examples improve success on unseen, reachable cube placements when example
count is controlled. H3 predicts that prior ViKi examples reduce the number
of simulation interactions needed to reach a given success rate. H4 predicts
that a camera-based policy trained from a successful state-based teacher can
perform the task when evaluated using only deployment-available observations.
Failure at any gate is a result to report, not a reason to relabel a geometric
plan as a physically successful demonstration.

\section{Contributions}
\label{sec:contributions}

The implemented contribution is an offline pipeline from calibrated RGB-D
recordings to hand trajectories, object tracks, and UR10 retarget plans. The
experimental contribution sought here is a reproducible assessment of
processing yield, simulated physical validity, object-relative variation,
and the marginal value of ViKi-derived data for state and visual policies.
Simulation replay, learning outcomes and real-world transfer are not claimed
until measured. The proposed experiment uses no motion-capture ground truth
or physical UR10; geometric consistency metrics must therefore not be
presented as absolute pose accuracy.

\section{Thesis Structure}
\label{sec:structure}

Chapter~\ref{ch:lit} reviews capture, retargeting and demonstration-driven
robot learning. Chapter~\ref{ch:method} defines the implemented pipeline and
the planned simulation study. Chapter~\ref{ch:impl} separates current evidence
from experiments still to run. Chapter~\ref{ch:disc} treats uncertainty and
validity, and Chapter~\ref{ch:concl} states the claims that the eventual
results can and cannot support.

\chapter{Literature Review}
\label{ch:lit}

\section{Preamble}
\label{sec:lit-preamble}

This chapter surveys the state of knowledge relevant to converting passively
recorded human manipulation demonstrations into robot-executable datasets. It
is organised not chronologically but by the architectural elements of the
proposed system, so that for every design decision taken in
Chapter~\ref{ch:method} the alternatives considered and the reason for the
choice are visible.

Section~\ref{sec:lit-hand} examines how the human hand is represented, from
full parametric meshes to the minimal wrist-pose-plus-gripper abstraction.
Section~\ref{sec:lit-retarget} contrasts sequential and single-cost-functional
retargeting formulations and reviews the solvers used to realise them.
Section~\ref{sec:lit-repr} compares object-centric and world-frame task
representations. Section~\ref{sec:lit-calib} covers multi-camera RGB-D
calibration, depth fusion and cross-device synchronisation, with emphasis on
reported quantitative accuracy. Section~\ref{sec:lit-dataset} reviews dataset
formats, downstream policies and demonstration curation.
Section~\ref{sec:positioning} positions \viki\ against the nearest published
systems, and Section~\ref{sec:lit-conclusion} states the resulting
methodological choices.

Sources were drawn from arXiv (\texttt{cs.RO}, \texttt{cs.CV}, \texttt{cs.LG}),
the proceedings of CoRL, RSS, ICRA, IROS, CVPR and NeurIPS, peer-reviewed
sensor and calibration literature, and vendor documentation for the depth
cameras and solvers employed. Where a source is a recent preprint whose
peer-review status could not be confirmed, this is stated explicitly.

\section{Representing the Human Hand}
\label{sec:lit-hand}

\subsection{Full parametric reconstruction}

A large body of work recovers the complete kinematics of the demonstrator.
OKAMI \cite{okami} combines SMPL-H for the body with MANO for the hand and
retargets the two separately. DexCanvas \cite{dexcanvas} builds a human--robot
correspondence dataset by fitting MANO to data from twenty-two infrared motion
capture cameras and two RGB-D sensors. DexMV \cite{dexmv} likewise recovers a
full MANO hand together with object pose from video before translating the
demonstration into simulation.

Full reconstruction yields a rich representation---the pose of every finger
joint---and is necessary when the target embodiment is a multi-fingered
dexterous hand. It also introduces a heavy perception stack and an additional,
poorly characterised error source: the mesh fit itself. For a parallel-jaw
gripper, most of that representational richness is discarded downstream.

\subsection{Minimal representation: wrist pose and gripper state}

The alternative, adopted in this work, represents the hand as a wrist pose in
$\SE$ together with a scalar or binary gripper state. This is not a concession
made for tractability; it is the representation used by the mainstream of
current teleoperation systems. Open-TeleVision \cite{opentv} maps the operator's
wrist pose to the arm end-effector through closed-loop inverse kinematics and
encodes a parallel gripper by the single thumb--index vector. Bunny-VisionPro
\cite{bunny} and ACE \cite{ace} adopt the same decomposition, ACE additionally
demonstrating that it transfers across dexterous hands, parallel grippers and
quadrupeds.

Methodological support for down-weighting the wrist comes from the ablation
study of Xin et al.\ \cite{keyobj}, who assemble a comprehensive retargeting
objective---global hand pose, overall hand shape, relative fingertip positions
and fingertip orientations---and remove terms one at a time. They find that
precise tracking of the global wrist pose is frequently unnecessary and can
degrade fingertip accuracy, and consequently replace the wrist-position
objective with a thumb-tip objective while regularising wrist orientation only
weakly. The implication for a wrist-plus-gripper representation is that the
wrist target should enter the optimisation as a weighted soft task rather than
a hard constraint.

\subsection{Egocentric and monocular capture}

Several systems reduce capture-hardware requirements further. EgoMimic
\cite{egomimic} records bare-handed demonstrations through Project Aria glasses
and co-trains on human and robot data. DexWild \cite{dexwild} collects
in-the-wild bare-hand data through a low-cost portable device and likewise
co-trains. A recent egocentric RGB-D system \cite{shadowing} generates
trajectories offline using analytic damped-least-squares inverse kinematics and
replays them on a low-cost arm. These works confirm a trend toward lighter
capture rigs, but each accepts monocular or single-view geometry, and none
combines it with the metric multi-view accuracy that the present work requires.

\section{Retargeting Formulations}
\label{sec:lit-retarget}

\subsection{The single cost functional}

The prevailing modern formulation poses retargeting as a constrained nonlinear
least-squares or quadratic program over the robot configuration $\bq_t$ at each
timestep, minimising a weighted sum of a data-fidelity term and a temporal
smoothness term subject to joint limits. AnyTeleop \cite{anyteleop} gives the
canonical instance: the objective sums, over keypoints, the squared residual
between a scaled human keypoint vector and the corresponding robot
forward-kinematics vector, plus a first-order penalty on configuration change
between consecutive frames. The accompanying \texttt{dex-retargeting} library
exposes a vector optimiser for real-time teleoperation, a position optimiser
explicitly recommended for offline dataset processing, and a contact-aware
DexPilot optimiser; it is built on Pinocchio and solved with sequential
quadratic programming.

Two refinements are directly relevant here. First, per-term confidence
weighting: in the formulation reported by Open-TeleVision \cite{opentv}, each
keypoint residual is scaled by a distance-dependent switching weight that
increases when a pinch is detected. The mechanism is precisely analogous to
weighting by a per-frame perception confidence, and it establishes that
weighting individual data terms by a scalar signal is accepted practice.
Second, second-order smoothness: DexFlow \cite{dexflow} augments the standard
objective with a dynamic smoothing weight and Hessian regularisation to enforce
$C^2$ continuity, penalising curvature and not merely velocity. Robustification
by a Huber loss on the data term, reported in recent dexterous teleoperation
work \cite{dexteleop0}, provides a mechanism for rejecting perception outliers
inside the optimiser rather than smoothing them away beforehand.

\subsection{The sequential pattern and its critique}

Against this, the sequential pattern---filter the observed task-space
trajectory, solve inverse kinematics per frame, filter the joint output---
remains common in offline video-to-robot pipelines. Its deficiency is
structural rather than empirical. Filtering is applied in task space, whereas
the smoothness requirement is on the joint trajectory; because inverse
kinematics is nonlinear, a smooth Cartesian input does not imply a smooth joint
output, particularly near singularities. Pre-filtering also attenuates genuine
fast motion indiscriminately with noise. In the single-functional formulation
the smoothness penalty acts directly on $\bq_t$, so joint-space smoothness is
obtained by construction.

\subsection{Offline pipelines and their retargeting operators}

ORION \cite{orion} does not perform per-frame hand retargeting: it constructs an
object-centric plan over sparse keyframes from an open-world object graph.
OKAMI \cite{okami} warps a reference trajectory to a new object position and
then solves inverse kinematics---the sequential pattern. HOWTransfer
\cite{howtransfer}, the closest published pipeline in workflow terms, performs
passive markerless capture and transfers trajectories using flow-matching
$\SE$ grasp sampling with Laplacian trajectory editing, not a single weighted
cost functional; its representation is hand-centric rather than object-centric.
No offline pipeline located in this review adopts the AnyTeleop-style single
functional.

\subsection{Solvers}

Pinocchio underlies the retargeting stack of AnyTeleop, Open-TeleVision and
ACE. Pink \cite{pink} exposes weighted inverse kinematics directly as a
quadratic program: each task contributes a quadratic residual with a weight
matrix normalising task coordinates to common units, while joint limits,
velocity limits and self-collision barriers enter as inequality constraints,
optionally with control-barrier formulations. This structure admits the entire
proposed cost functional natively---a frame task for the wrist target, posture
and velocity regularisation tasks for smoothness, and limits as constraints---
solved as one program per frame. The solver choice made in
Chapter~\ref{ch:method} therefore lies on the mainstream path rather than
diverging from it.

\section{Task Representation}
\label{sec:lit-repr}

Object-centric representations dominate recent work on cross-environment and
cross-embodiment generalisation. ORION \cite{orion} represents a task as the
three-dimensional object trajectory relative to initial and final object
positions, arguing that this generalises across visual backgrounds, camera
shifts, spatial layouts and novel object instances, and that full-body
reconstruction of the demonstrator is unnecessary. SPOT \cite{spot} encodes each
task as the $\SE$ object-pose trajectory relative to the target explicitly in
order to decouple embodiment-specific actions from sensory input, enabling
learning from action-free human demonstrations, and conditions a diffusion
policy on the object trajectory. Flow-based formulations \cite{novaflow} make an
equivalent argument while relaxing the rigid-object assumption that
pose-trajectory methods require.

OKAMI \cite{okami} occupies an intermediate position: the representation is
demonstrator-centric, but object-aware warping adapts trajectories to new object
positions---an implicit acknowledgement that pure world-frame replay fails when
the object moves. Vi-RoIL \cite{viroil} is object-centric by construction,
expressing wrist trajectories and keyframe hand poses in the object frame, but
recovers metric scale from a monocular view of a marker cube and evaluates only
in simulation.

The trade-off is that object-centric representations require robust object pose
or segmentation, and pose-trajectory variants assume rigidity, whereas
world-frame trajectories are simpler but tied to the recording layout. Storing
both, as proposed in Section~\ref{sec:objcentric}, defers the choice to the
downstream consumer at negligible storage cost.

\section{Multi-Camera RGB-D Calibration, Fusion and Synchronisation}
\label{sec:lit-calib}

HO-Cap provides eight calibrated RealSense views and semi-automatic MANO
hand--object annotations \cite{hocap}. It is an external reference for the
hand-joint stage, not an optical-motion-capture measurement of our two-Kinect
rig. Camera selection on the evaluated sequences must be reported as pilot
selection rather than held-out performance.

\subsection{Fiducial choice}

ChArUco boards combine the sub-pixel corner accuracy of a checkerboard with the
robust identification and occlusion tolerance of ArUco markers. A $6\times8$
board provides up to thirty-five corners per view against four for a single
AprilTag, and remains detectable under partial occlusion and at the image
boundary where a plain checkerboard fails. Reported reprojection errors for
ChArUco and ArUco calibration are typically below half a pixel. AprilTags are
easier to detect at long range and oblique incidence and suffice for
extrinsic-only estimation, but are ill-suited to intrinsic calibration owing to
the small number of corners. Current practice is therefore to use ChArUco for
intrinsics and either fiducial for extrinsics, preferring AprilTags for oblique
overhead views.

\subsection{Reported accuracy of multi-Azure-Kinect rigs}

Romeo et al.\ \cite{romeo} provide the cleanest direct comparison of calibration
procedures for multi-Kinect rigs, reporting that three-dimensional procedures
outperform two-dimensional chessboard approaches, with mean distances between
corresponding point-cloud points of $21.4$~mm at colour resolution and
$9.9$~mm at infrared resolution for a static scene, and $20.9$~mm and $7.4$~mm
respectively for a dynamic scene; body-joint alignment error was $35.4$~mm.
Darwish et al.\ \cite{darwish} calibrate five Azure Kinects against a
ChArUco-coded cube, combining reprojection error with point-to-plane and
patch-to-plane terms to compensate depth noise. Kim et al.\ \cite{kimcal}
perform coarse registration from body tracking followed by feature-match
refinement, retaining above ninety-five per cent inliers. Bu and Lee
\cite{bulee} use markerless coarse calibration with coloured ICP refinement and
report---importantly for rig geometry---that ICP fails for near-opposing camera
placements owing to insufficient point-cloud overlap.

Single-sensor specifications bound the achievable accuracy from below. The
Azure Kinect DK is specified at a random error standard deviation of at most
$17$~mm and a typical systematic error below $11$~mm plus $0.1$ per cent of
distance \cite{mskinect}; independent evaluation confirms these figures
\cite{tolgyessy}, and comparative work finds improved spatial and temporal
accuracy relative to Kinect~v2 beyond $2.5$~m \cite{kurillo}. The Intel
RealSense D435i is specified at depth accuracy below one per cent of distance
\cite{intel}. The practical consequence is that centimetre-level, not
millimetre-level, agreement should be expected between calibrated views.

\subsection{Depth fusion}

Naive averaging of overlapping depth observations is inferior to
confidence-weighted fusion. Weighting by distance and by the angle between the
surface normal and the viewing direction---so that surfaces observed near
perpendicular contribute more than obliquely observed ones---is established
practice \cite{fusionpatent}. Learned confidence maps estimated from the depth
signal itself are used to re-weight depth features before fusion
\cite{zeng}, and surfel-based reconstruction attaches an explicit confidence
value to each element and performs curvature-weighted registration
\cite{hrbf}. Both sensor families used in this work expose per-pixel
invalidation or confidence signals that can drive such weights directly.

\subsection{Cross-device synchronisation}

Azure Kinect devices support hardware synchronisation through daisy-chained or
star-wired trigger cabling, with subordinate devices required to start before
the master and depth exposures staggered to prevent laser interference; the
vendor documentation states explicitly that equal timestamps between depth and
colour streams, or between cameras, are not guaranteed \cite{mskinect}.
RealSense devices provide an analogous inter-camera synchronisation mode, with
the counter-intuitive property that timestamps appear aligned when hardware
synchronisation is disabled and exhibit bounded drift when it is enabled, the
drift itself serving as verification \cite{intel, ridgerun}.

The pitfall for a heterogeneous rig is that hardware clocks on different device
families are not mutually comparable, since each counts its own cycles from its
own epoch. Documented practice for such rigs is to stamp every arriving frame
with a single host-side monotonic clock and to align streams by nearest
timestamp in that common domain, accepting millisecond-scale precision
\cite{ridgerun}. Because Azure Kinect and RealSense triggers are mutually
incompatible, this is the only available mechanism for aligning the two
families, and it is the mechanism adopted in Section~\ref{sec:sync}.

\subsection{Hand--eye calibration}

The OpenCV \texttt{calibrateHandEye} routine implements the Tsai--Lenz,
Park--Martin, Horaud--Dornaika, Andreff and Daniilidis solvers. A comparative
noise study \cite{alihandeye} finds the Daniilidis, Chou and Park methods most
robust to translation noise, with Daniilidis yielding the best translation
estimates, while Tsai, Li and Lu degrade as rotation noise grows. Methods that
decouple rotation from translation are faster but propagate rotation error into
the translation estimate; Daniilidis solves both simultaneously through dual
quaternions. Because the present system anchors cameras and robot to a shared
ChArUco workspace frame, classical hand--eye estimation is largely bypassed;
where an explicit camera-to-robot transform is required, Daniilidis is used.

\section{Dataset Generation and Downstream Policies}
\label{sec:lit-dataset}

LeRobot-compatible datasets together with ACT and Diffusion Policy are
potential downstream consumers. The implemented default export is a NumPy
trajectory bundle, not a replay-validated LeRobot dataset. OpenVLA-OFT \cite{openvlaoft} reports that parallel decoding, action
chunking, a continuous action representation and $L1$ regression raise average
success across four task suites from $76.5$ to $97.1$ per cent while increasing
action-generation throughput twenty-six-fold, and is therefore a credible
alternative consumer of the generated data.

Curation matters as much as generation. Demo-SCORE \cite{demoscore} trains a
classifier on policy rollouts to filter low-quality demonstrations and reports
an absolute success-rate improvement exceeding fifteen to thirty-five per cent
over training on all demonstrations. Smoothness-based scoring of trajectories
has been reported to yield comparable gains at a fraction of the data
\ph{[citation to be confirmed]}. This motivates treating dataset curation as an
explicit pipeline stage rather than an afterthought, and---in the present
design---eventually screening physical replay before export. This screening
is not implemented in the current system.

\section{Positioning of the Proposed System}
\label{sec:positioning}

Table~\ref{tab:positioning} compares \viki\ against the nearest published
systems along the six criteria that jointly define the proposed architecture:
(C1) passive markerless capture without teleoperation hardware; (C2)
multi-camera metrically calibrated RGB-D sensing; (C3) object-centric dense
trajectory representation; (C4) retargeting as a single confidence-weighted
cost functional; (C5) export of an imitation-learning-ready dataset; and (C6)
physical replay used to validate the dataset before export.

\begin{table}[htbp]
\caption{Comparison of the Proposed System Against the Nearest Published Work\\
Along the Six Criteria Defined in Section~\ref{sec:positioning}}
\label{tab:positioning}
\centering
\small
\begin{tabular}{lcccccc}
\toprule
System & C1 & C2 & C3 & C4 & C5 & C6 \\
\midrule
\viki\ (proposed)      & \checkmark & \checkmark & \checkmark & \checkmark & \checkmark & \checkmark \\
HOWTransfer \cite{howtransfer} & \checkmark & $\sim$ & --- & --- & $\sim$ & \checkmark \\
Vi-RoIL \cite{viroil}          & $\sim$ & --- & \checkmark & --- & $\sim$ & --- \\
ORION \cite{orion}             & \checkmark & $\sim$ & \checkmark & --- & --- & --- \\
OKAMI \cite{okami}             & \checkmark & $\sim$ & $\sim$ & --- & $\sim$ & --- \\
EgoMimic \cite{egomimic}       & \checkmark & --- & --- & --- & $\sim$ & --- \\
DexWild \cite{dexwild}         & $\sim$ & --- & --- & --- & $\sim$ & --- \\
Tool-as-Interface \cite{toolif}& \checkmark & $\sim$ & $\sim$ & --- & $\sim$ & --- \\
RoboPaint \cite{robopaint}     & --- & \checkmark & $\sim$ & --- & \checkmark & \checkmark \\
DexCanvas \cite{dexcanvas}     & --- & \checkmark & --- & --- & $\sim$ & --- \\
\bottomrule
\end{tabular}

\vspace{1mm}
{\footnotesize \checkmark~= satisfied; $\sim$~= partially satisfied;
---~= not satisfied.}
\end{table}

No located system satisfies all six criteria simultaneously. HOWTransfer
\cite{howtransfer} is closest in workflow, satisfying C1 and C6, but is
hand-centric rather than object-centric, uses calibrated stereo RGB rather than
metric RGB-D, and does not employ a single cost functional. Vi-RoIL
\cite{viroil} is closest in representation, satisfying C3, but relies on a
monocular view of a marker cube and is evaluated only in simulation. The
intersection of the online joint-optimisation retargeting literature with the
offline object-centric dataset-generation literature is therefore unoccupied,
and it is that intersection which this work targets.

Several of the systems listed are recent preprints
\cite{howtransfer, robopaint, shadowing, dexcanvas} whose peer-review status
could not be confirmed at the time of writing; their reported results are
treated as provisional.

\section{Conclusion of the Review}
\label{sec:lit-conclusion}

The review supports the following methodological choices, which
Chapter~\ref{ch:method} implements.

\begin{enumerate}[leftmargin=2em]
  \item \emph{Hand representation.} A wrist pose in $\SE$ with a continuous
        gripper opening is adopted, following \cite{opentv, bunny, ace}, and the wrist
        target enters as a weighted soft task following the ablation evidence
        of \cite{keyobj}. Parametric mesh reconstruction is rejected as
        unnecessary for a parallel-jaw target and as an additional error source.
  \item \emph{Retargeting.} A single weighted cost functional is adopted,
        following \cite{anyteleop, dexflow}, with confidence weighting after
        \cite{opentv} and Huber robustification after \cite{dexteleop0}. The
        sequential smooth--IK--smooth pattern is rejected on the structural
        grounds set out in Section~\ref{sec:lit-retarget}.
  \item \emph{Solver.} Pink on Pinocchio \cite{pink} is adopted, since it
        expresses the entire functional natively as a constrained quadratic
        program and is the solver family used by the reference teleoperation
        systems.
  \item \emph{Task representation.} Object-relative trajectories are adopted
        following \cite{orion, spot}, and stored alongside workspace-anchored
        trajectories so that the choice is deferred downstream.
  \item \emph{Sensing.} ChArUco workspace anchoring is adopted; centimetre-level
        cross-camera agreement is assumed on the evidence of \cite{romeo};
        camera placement preserves substantial view overlap on the evidence of
        \cite{bulee}; depth is fused with distance, incidence-angle and sensor
        confidence weights following \cite{fusionpatent, zeng}.
  \item \emph{Synchronisation.} The two Azure Kinects support hardware
        synchronisation, and the application uses a host-side monotonic
        clock for cross-device frame pairing \cite{ridgerun}. Local timing
        accuracy cannot be re-evaluated without the source recordings.
  \item \emph{Curation.} Physical replay screening is a proposed validation
        stage, motivated by \cite{demoscore} and \cite{howtransfer}; the
        current implementation does not perform it.
\end{enumerate}

\chapter{Design and Methodology}
\label{ch:method}

\section{System Overview}
\label{sec:overview}

\viki\ is an offline pipeline: demonstrations are recorded, processed and
exported in distinct stages, with no requirement that any stage run at
interactive rates. This is a deliberate design choice. Relaxing the real-time
constraint that shapes teleoperation systems permits offline reconstruction
and robot-plan inspection. Physical replay is a planned validation interface,
not an implemented or evaluated stage.

The pipeline comprises six stages, shown in Fig.~\ref{fig:pipeline}.

\begin{enumerate}[leftmargin=2em]
  \item \emph{Capture.} Synchronised RGB-D streams are recorded from a
        calibrated multi-camera rig while a demonstrator performs the task
        bare-handed.
  \item \emph{Perception.} Hand landmarks are detected per view, lifted to
        three dimensions using sensor depth, and fused across views with
        confidence weights to yield a wrist pose in $\SE$ and a gripper state.
  \item \emph{Representation.} The wrist trajectory is expressed both relative
        to the manipulated object and relative to the ChArUco workspace anchor.
  \item \emph{Retargeting.} A robot joint trajectory is obtained by minimising
        a single weighted cost functional subject to kinematic and collision
        constraints.
  \item \emph{Replay validation (stub).} The interface exists, but no physical
        robot execution or replay-based acceptance is implemented.
  \item \emph{Export.} A NumPy trajectory bundle and JSON manifest can be
        written without implying physical validity; LeRobot export is optional
        and subject to additional requirements.
\end{enumerate}

\begin{figure}[htbp]
\centering
\begin{tikzpicture}[
  node distance=4mm,
  box/.style={draw, rounded corners=1pt, align=center, font=\footnotesize,
              minimum height=8mm, text width=22mm, inner sep=2pt},
  arr/.style={-{Latex[length=2mm]}, thick}]
\node[box] (cap) {Capture\\\scriptsize 2$\times$Kinect};
\node[box, right=of cap] (per) {Perception\\\scriptsize landmarks + fusion};
\node[box, right=of per] (rep) {Representation\\\scriptsize object-relative $\SE$};
\node[box, below=8mm of rep] (ret) {Retargeting\\\scriptsize single cost functional};
\node[box, left=of ret] (val) {Replay\\\scriptsize stub};
\node[box, left=of val] (exp) {Export\\\scriptsize NumPy bundle};
\draw[arr] (cap) -- (per);
\draw[arr] (per) -- (rep);
\draw[arr] (rep) -- (ret);
\draw[arr] (ret) -- (val);
\draw[arr] (val) -- (exp);

\end{tikzpicture}
\caption{Stages of the \viki\ pipeline. Capture, perception and representation
are embodiment-agnostic; retargeting is specific to the target manipulator.
Replay is a stub, and export does not certify physical execution.}
\label{fig:pipeline}
\end{figure}

\section{Hardware and Workspace}
\label{sec:hardware}

The local capture setup uses two Azure Kinect DK cameras observing a
tabletop workspace. UR10 with a parallel-jaw gripper is the intended
simulation target, not a physical manipulator mounted in this setup.
Reachability of the unavailable local demonstrations cannot be quantified.
The Azure Kinect time-of-flight depth characteristics are documented by the
vendor \cite{mskinect} and studied independently \cite{tolgyessy}; neither
source validates the current rig's calibration or fused trajectory accuracy.

The two Kinects are placed to preserve substantial overlap in their fields of
view. This is a direct consequence of the finding of Bu and Lee \cite{bulee}
that iterative-closest-point refinement fails between near-opposing views for
lack of point-cloud overlap; near-orthogonal placement retains both
triangulation baseline and overlap.

\section{Calibration}
\label{sec:calib}

\subsection{Intrinsics and extrinsics}

Intrinsic parameters for every colour and depth stream are estimated from
ChArUco board observations. Extrinsic calibration follows the workspace-anchor
principle: rather than expressing cameras relative to one another, each camera
is calibrated against a ChArUco board rigidly fixed to the workspace, defining
a common world frame $W$. For camera $k$ this yields
\begin{equation}
  \bT^{W}_{C_k} \in \SE,
  \label{eq:extrinsic}
\end{equation}
the transform from the camera frame $C_k$ to the workspace frame $W$. Anchoring
to the workspace rather than to a reference camera has two consequences: adding
or moving a camera does not invalidate the calibration of the others, and the
robot base can be registered into the same frame without a separate
camera-to-camera chain.

\subsection{Robot registration}

The manipulator base frame $R$ is registered to $W$ by observing a ChArUco
target held in a known end-effector pose across a set of configurations, giving
$\bT^{W}_{R}$ directly. Where an explicit camera-to-robot transform is required
instead, the Daniilidis dual-quaternion solver is used, following the
robustness evidence of \cite{alihandeye}, with fifteen to twenty poses spanning
a hemisphere.

\section{Synchronisation}
\label{sec:sync}

The two Azure Kinects are hardware-synchronised in master/subordinate
configuration, with subordinate devices started before the master and depth
exposures staggered to avoid mutual laser interference \cite{mskinect}. The
RealSense cannot share that trigger; its synchronisation mode is incompatible
at the electrical level.

Cross-family alignment therefore uses a single host-side monotonic clock. Every
frame, from every device, is stamped on arrival with
\begin{equation}
  \tau = \texttt{time.monotonic\_ns()},
  \label{eq:hostclock}
\end{equation}
and streams are aligned by nearest neighbour in $\tau$. Two points deserve
emphasis. First, per-device hardware timestamps are \emph{not} used for
cross-device arithmetic: distinct device families count their own cycles from
their own epochs, and the vendor documentation states explicitly that equal
timestamps across cameras are not guaranteed \cite{mskinect}. Second, a
monotonic clock rather than wall-clock time is required, since wall-clock time
may be adjusted discontinuously by time synchronisation daemons, which would
corrupt interval computations mid-recording.

A future local evaluation must preserve paired device and host timestamps to
measure alignment jitter and drift. No local synchronisation statistic is
reported here; the HO-Cap benchmark cannot validate the Kinect rig.

\section{Perception and Multi-View Reconstruction}
\label{sec:fusion}

\subsection{Landmark detection and lifting}

Hand landmarks are detected in each colour stream, yielding for view $k$ at
time $t$ a set of image-plane landmarks with associated visibility scores
$v^{k}_{i,t} \in [0,1]$. Each landmark is lifted to three dimensions using the
registered depth map rather than a learned depth estimate, exploiting the
time-of-flight measurement directly, and transformed into the workspace frame
by \eqref{eq:extrinsic}.

\subsection{Multi-view triangulation}

When a landmark is seen by two or more calibrated views, its position is not
averaged from per-view depth lifts but estimated directly from the image
observations, with depth entering as an additional soft residual rather than as
the anchor. Averaging per-view lifts inherits the depth error of every view;
triangulation is constrained by the camera geometry, which is known to
calibration accuracy and does not degrade with range.

For landmark $i$ at time $t$, let $\mathcal{K}_{i,t}$ be the views whose
detector score clears a threshold and whose observation is finite. Every pair
$(a,b) \subset \mathcal{K}_{i,t}$ yields a linear direct-linear-transform
estimate, which is retained only if it lies in front of both cameras and if the
two viewing rays subtend at least a minimum angle $\theta_{\min}$; a pair with a
short effective baseline is ill-conditioned regardless of how confident the
detector is. Each surviving hypothesis is scored against all views in
$\mathcal{K}_{i,t}$ by the number of reprojection inliers, and the best is
refined over its inlier set by
\begin{equation}
  \hat{\bx}_{i,t} \;=\;
  \argmin_{\bx}\;
  \sum_{k \in \mathcal{K}^{\text{in}}_{i,t}}
  \rho\!\left(\norm{\pi_k(\bx) - \bu^{k}_{i,t}}^{2}\right)
  \;+\;
  \lambda_d \sum_{k} \beta^{k}_{i,t}\,
  \rho_d\!\left(\left(z_k(\bx) - d^{k}_{i,t} - \delta_i\right)^{2}\right),
  \label{eq:triangulate}
\end{equation}
where $\pi_k$ is the pinhole projection of view $k$, $\bu^{k}_{i,t}$ the
observed pixel, $z_k$ the depth of $\bx$ in that camera, $d^{k}_{i,t}$ the
measured depth, and $\delta_i$ a per-landmark offset from the skin surface to
the joint centre---depth sees the skin, the joint sits inside the finger. The
depth weight $\beta^{k}_{i,t}$ is attenuated by the local depth spread and set
to zero for invalid samples, so that a single bad silhouette pixel cannot move
an otherwise well-constrained point. Observations are undistorted once, before
the linear estimate, and every projection downstream is a plain pinhole map;
distortion is never re-applied inside the residual.

The estimate carries a transparent geometric quality
\begin{equation}
  q_{i,t} \;=\;
  \underbrace{\frac{\left|\mathcal{K}^{\text{in}}_{i,t}\right|}
                   {\left|\mathcal{K}_{i,t}\right|}}_{\text{agreement}}
  \cdot
  \underbrace{\operatorname{clip}\!\left(
      1 - \frac{\bar{e}_{i,t}}{2 e_{\text{in}}},\,0,\,1\right)}_{\text{reprojection}}
  \cdot
  \underbrace{\operatorname{clip}\!\left(
      \frac{\theta_{i,t}}{\theta_{\text{ref}}},\,0,\,1\right)}_{\text{baseline}},
  \label{eq:quality}
\end{equation}
with $\bar{e}_{i,t}$ the mean reprojection error over inliers, $\theta_{i,t}$
the widest ray angle, and $e_{\text{in}}$ the inlier radius. This is a
transparent heuristic, not a calibrated variance, and it is used as such: it
weights evidence, and is never interpreted as an inverse covariance.

A property of \eqref{eq:quality} deserves explicit statement, because it
governs how much of a recording survives. The same $e_{\text{in}}$ is the
acceptance threshold, the scale of the robust loss $\rho$ in
\eqref{eq:triangulate}, and the denominator of the reprojection factor. The
factor is therefore designed to grade confidence out to $2 e_{\text{in}}$,
while the threshold admits nothing beyond $e_{\text{in}}$: half of the designed
grading range is unreachable by construction. Widening $e_{\text{in}}$ does not
introduce a soft-weighting mechanism, it permits the existing one to operate
over its own range. A consequence is that $q$ is not comparable across settings
of $e_{\text{in}}$: the same residual scores differently under a different
denominator.

\subsection{Single-view fallback}

Where only one view observes a landmark, triangulation is undefined and the
estimate falls back to the depth lift of that view. The historical multi-view
path also combined per-view lifts by confidence-weighted averaging,
\begin{equation}
  \hat{\bx}_{i,t} \;=\;
  \frac{\sum_{k} w^{k}_{i,t}\, \bx^{k}_{i,t}}
       {\sum_{k} w^{k}_{i,t}},
  \qquad
  w^{k}_{i,t} \;=\;
  \underbrace{v^{k}_{i,t}}_{\text{detector}}
  \cdot
  \underbrace{c^{k}_{i,t}}_{\text{sensor}}
  \cdot
  \underbrace{\frac{1}{\left(d^{k}_{i,t}\right)^{2}}}_{\text{range}}
  \cdot
  \underbrace{\max\!\left(0, \cos\theta^{k}_{i,t}\right)}_{\text{incidence}},
  \label{eq:fusion}
\end{equation}
where $v^{k}_{i,t}$ is the detector score, $c^{k}_{i,t}$ the sensor depth
validity at that pixel, $d^{k}_{i,t}$ the range, and $\theta^{k}_{i,t}$ the
incidence angle. The range and incidence factors follow established practice in
geometric depth fusion, where surfaces observed near perpendicular are weighted
above obliquely observed ones because they are measured more accurately
\cite{fusionpatent}; the detector and sensor factors follow confidence-map
re-weighting \cite{zeng}. The product form ensures that a zero in any
factor---an undetected landmark, an invalid depth pixel, a grazing view---
removes that view's contribution entirely rather than attenuating it.

This path is retained for reproducibility of earlier recordings and for
single-camera operation. It is not the path used for the results reported here,
because averaging per-view lifts propagates depth error into the fused point,
whereas \eqref{eq:triangulate} constrains the point by geometry and admits
depth only as evidence.

\subsection{Wrist pose and gripper state}

The wrist pose $\bT^{W}_{H,t} \in \SE$ is obtained by fitting an orthonormal
frame to the fused wrist and knuckle landmarks. The gripper state is the thumb--index
separation, mapped to a normalised opening,
\begin{equation}
  g_t \;=\;
  \operatorname{clip}\!\left(
    \frac{\norm{\hat{\bx}_{\text{thumb},t} - \hat{\bx}_{\text{index},t}}
          - \tau_{\text{closed}}}
         {\tau_{\text{open}} - \tau_{\text{closed}}},\,0,\,1\right),
  \label{eq:gripper}
\end{equation}
temporally filtered to suppress chatter. An earlier formulation thresholded the
same distance into a binary open/closed indicator; it was replaced because a
parallel-jaw gripper is a continuously commandable actuator, and quantising the
one degree of freedom that carries the grasp discards precisely the approach
and release phases that distinguish a successful demonstration from a failed
one. This single scalar remains the entire hand representation beyond the wrist
frame, following the convention established for parallel-jaw targets in
\cite{opentv, ace}.

\section{External Hand-Joint Reference Protocol}
\label{sec:hocap-method}

The main external cohort contains all 28 selected single-hand HO-Cap
pick-and-place sequences (24 right-hand and four left-hand, 29,706 frames)
from the pinned HO-Cap mirror revision recorded in the archived batch
status. All eight RealSense RGB-D views are processed with MediaPipe observations,
calibrated triangulation and the frozen \texttt{hocap-mediapipe-triangulated-v1}
profile. The HoloLens stream is excluded. These sequences are separate from
the local Kinect cube cohort and cannot validate its calibration. Published
semi-automatic MANO 21-joint positions are transformed into the
\texttt{hocap\_tag\_1} frame. Extract and Prepare do not use the reference
hand labels. Predictions are paired by exact frame index and synthetic
capture timestamp, without root alignment, Procrustes fit, temporal search
or reference-guided filtering. MPJPE is the mean Euclidean error over finite
matched joints; coverage is matched joints divided by valid reference joints.
We report PCK@20/50~mm and the per-sequence 95th percentile alongside MPJPE.
The aggregate MPJPE and PCK are weighted by matched joint counts; an
aggregate P95 cannot be recovered from per-sequence quantiles. Incomplete
Prepare stages are excluded.

For each of the four named objects in a sequence, a separate tracker receives
\emph{published instance masks} and depth from the eight views. It does not
predict segmentation. Model confidence must reach 0.35. A single
right-multiplied object-local SE(3) transform is fit at the first confident
frame using the published object pose. Later matched frames are scored by
translation distance and geodesic SO(3) angle in the same world frame;
coverage is tracked frames divided by valid reference frames. Thus these are
\emph{relative tracking-drift} scores conditional on oracle masks and a
reference-assisted first-frame gauge choice, not absolute 6-DoF localization.
No symmetry-equivalent orientation metric or moving-only subset is applied.
The 28-scene selection is a task-and-handedness subset, not an official
held-out HO-Cap benchmark split. Earlier two-sequence, 120-frame two-versus-
eight-view results remain a development pilot: its camera pair was selected
on the same frames and cannot be treated as held-out evidence.

\section{Object-Centric Representation}
\label{sec:objcentric}

Let $\bT^{W}_{O,t}$ denote the pose of the manipulated object in the workspace
frame. Each demonstration is stored in two forms. The workspace-anchored form
is $\bT^{W}_{H,t}$ directly; the object-relative form is
\begin{equation}
  \bT^{O}_{H,t} \;=\; \left(\bT^{W}_{O,t}\right)^{-1} \bT^{W}_{H,t}.
  \label{eq:objrel}
\end{equation}

The object-relative form is what transfers. If the same task is to be performed
with the object at a new pose $\bT^{R}_{O'}$ expressed in the robot frame, the
desired end-effector pose at time $t$ is
\begin{equation}
  \bT^{R,\mathrm{des}}_{E,t} \;=\;
  \bT^{R}_{O'} \; \bT^{O}_{H,t} \; \bT^{H}_{E},
  \label{eq:transfer}
\end{equation}
where $\bT^{H}_{E}$ is a fixed offset relating the human wrist frame to the
robot end-effector frame, determined once per embodiment. Trajectories are
retained densely rather than reduced to keyframes, since the continuous
velocity profile of the demonstration carries task-relevant information---
approach speed, contact timing---that sparse keyframes discard.

Both representations are exported. Storing both costs little and defers to the
downstream consumer a choice that the literature does not settle universally:
object-relative trajectories generalise across object placement
\cite{orion, spot} but depend on object pose estimation, whereas
workspace-anchored trajectories are simpler and adequate within a fixed layout.

\section{Retargeting as a Single Cost Functional}
\label{sec:retarget}

\subsection{Formulation}

Let $\bq_t \in \R^{n}$ be the manipulator configuration at time $t$ and
$f(\bq_t) \in \SE$ the forward kinematics of the end-effector. The task-space
residual with respect to the desired pose \eqref{eq:transfer} is the
logarithmic map of the pose error,
\begin{equation}
  \be_t(\bq_t) \;=\;
  \Log\!\left(
    \left(\bT^{R,\mathrm{des}}_{E,t}\right)^{-1} f(\bq_t)
  \right)^{\vee} \in \R^{6}.
  \label{eq:residual}
\end{equation}

Retargeting is then posed as the single constrained problem
\begin{equation}
\begin{aligned}
  \min_{\bq_{1:T}} \;\; J(\bq_{1:T}) \;=\; \sum_{t=1}^{T} \Big[\;
   & \underbrace{\omega_t \,\rho_\delta\!\left(
       \norm{\be_t(\bq_t)}_{\bm{\Sigma}}\right)}_{\text{data fidelity}}
   \;+\;
     \underbrace{\lambda_v \norm{\bq_t - \bq_{t-1}}^{2}}_{\text{velocity}}
   \\[1mm]
   &+\;
     \underbrace{\lambda_a \norm{\bq_t - 2\bq_{t-1} + \bq_{t-2}}^{2}}
       _{\text{acceleration}}
   \;+\;
     \underbrace{\lambda_r \norm{\bq_t - \bq_{\mathrm{ref}}}^{2}}
       _{\text{posture}}
  \;\Big]
\end{aligned}
\label{eq:cost}
\end{equation}
subject to
\begin{align}
  \bq_{\min} \;\le\; &\bq_t \;\le\; \bq_{\max}, \label{eq:cjoint}\\
  \norm{\bq_t - \bq_{t-1}} \;&\le\; \dot{\bq}_{\max}\,\Delta t,
    \label{eq:cvel}\\
  h_j(\bq_t) \;&\ge\; 0, \quad j = 1,\dots,m, \label{eq:ccoll}
\end{align}
where

\begin{center}
\begin{tabular}{@{\hspace{10mm}}p{28mm}@{}p{85mm}}
$\omega_t$ & per-frame perception confidence weight, defined in \eqref{eq:conf};\\
$\rho_\delta$ & Huber loss with transition parameter $\delta$;\\
$\bm{\Sigma}$ & diagonal matrix normalising translational and rotational residual units;\\
$\lambda_v, \lambda_a, \lambda_r$ & weights on velocity, acceleration and posture regularisation;\\
$h_j$ & collision and self-collision barrier functions.
\end{tabular}
\end{center}

\subsection{Confidence weighting}

The per-frame weight aggregates the geometric qualities \eqref{eq:quality} of
the landmarks that define the palm frame at time $t$,
\begin{equation}
  \omega_t \;=\;
  \left(\frac{1}{|\mathcal{I}|}\sum_{i \in \mathcal{I}} c_{i,t}\right)^{\alpha}
  \cdot \mathbb{1}\!\left[\text{frame valid}\right],
  \qquad
  c_{i,t} =
  \begin{cases}
    q_{i,t}, & \text{landmark triangulated,}\\
    0, & \text{landmark interpolated,}
  \end{cases}
  \label{eq:conf}
\end{equation}
with $\mathcal{I}$ the four landmarks defining the palm frame---wrist and the
index, middle and little metacarpophalangeal joints---and $\alpha$ controlling
how sharply low-confidence frames are discounted.

Equation \eqref{eq:conf} as written admits a frame whose palm was only
partially observed: a mean over four terms is non-zero whenever one of them is,
so a palm with a single triangulated landmark still contributes at roughly a
quarter weight, and the orientation it contributes was fitted partly through
interpolated points. A stricter diagnostic variant zeroes the frame unless every landmark in
$\mathcal{I}$ was triangulated,
\begin{equation}
  \omega^{\text{strict}}_t \;=\;
  \omega_t \cdot
  \mathbb{1}\!\left[\,\forall i \in \mathcal{I}:\; \text{$i$ triangulated}\,\right].
  \label{eq:confstrict}
\end{equation}
This variant is a proposed diagnostic, not a measured improvement in the
available evidence. It trades observation support for stricter provenance;
a retained local cohort is needed to quantify the effect on tracking and
coverage. The construction is analogous to distance-dependent switching
weights in online teleoperation retargeting \cite{opentv}: poorly observed
frames should exert less influence than well-observed frames.

\subsection{Robustification}

The Huber loss $\rho_\delta$ bounds the influence of gross outliers---a
mis-detected hand, a depth dropout that survives the confidence weighting---at
the level of the data term itself, rather than allowing a filtering stage to
smear the outlier across neighbouring frames \cite{dexteleop0}. This is the
mechanism that replaces pre-smoothing: outliers are down-weighted where they
enter, not attenuated after they have already propagated.

\subsection{Why a single functional}

Equation~\eqref{eq:cost} differs from the sequential alternative in one
structurally important respect. The smoothness penalties act on $\bq_t$, in the
space where smoothness is required, and they are traded against data fidelity
by the solver rather than applied blindly beforehand. In the sequential
pipeline, task-space filtering cannot guarantee joint-space smoothness because
the forward kinematics is nonlinear, and it cannot distinguish task-relevant
fast motion from noise because it operates without reference to the data term.
Here the trade-off is explicit: on high-confidence frames $\omega_t$ is large
and the solution tracks the observation; on low-confidence frames the
smoothness terms dominate and the solution interpolates.

Second-order smoothness, expressed by the acceleration term, follows the
$C^2$-continuity argument of DexFlow \cite{dexflow}. A trajectory that is
first-order smooth may still exhibit curvature discontinuities that appear as
jerk on the physical manipulator and as label noise to a policy trained on the
data.

\subsection{Solution}

The problem is solved with Pink on Pinocchio \cite{pink}. Each term of
\eqref{eq:cost} maps onto a weighted task in the quadratic program---a frame
task for the data term, posture and velocity regularisation tasks for the
smoothness terms---while \eqref{eq:cjoint}--\eqref{eq:ccoll} enter as
inequality constraints, with collision avoidance expressed through control
barrier functions. Levenberg--Marquardt damping stabilises the solution near
singularities.

Because the setting is offline, the problem is solved over the whole trajectory
rather than frame by frame, so the acceleration term is well defined at every
interior timestep and the solution is not constrained by the causality
requirement that an online solver faces. The trajectory is initialised from the
previous solution warm-started along $t$, and boundary conditions at $t=1,2$
are handled by replicating the initial configuration.

\subsection{Interpolation}

Where gaps remain after fusion---landmarks lost across all views---the
trajectory is completed by cubic spline interpolation over $\SE$ prior to
optimisation. Dynamic filters that impose a motion model, such as Kalman or
extended Kalman filters, are deliberately not applied: they introduce a prior
about the dynamics that is subsequently re-estimated by the optimiser, and the
interaction between the two priors is difficult to characterise. The optimiser
already imposes the required regularity through
$\lambda_v$ and $\lambda_a$.

\section{Replay Validation}
\label{sec:replay}

Physical replay remains an interface stub: no robot execution, measured
proprioception, contact outcome, or controller tracking residual is available.
Kinematic plans must not be presented as replay-validated demonstrations.
Simulator or physical validation would require recorded execution, contact
checks, explicit success criteria, and rejection of failed episodes; these
results are not part of the present evidence.

\section{Dataset Export}
\label{sec:export}

The implemented default export writes a self-contained NumPy trajectory bundle
and a JSON manifest. It does not imply physical replay or outcome screening.
An optional LeRobot-format export has additional eligibility requirements and
dependencies; neither format makes an unvalidated plan a physical
demonstration.

\section{Assumptions and Simplifications}
\label{sec:assumptions}

The following assumptions delimit the scope of the work and are revisited in
Section~\ref{sec:implications}.

\begin{enumerate}[leftmargin=2em]
  \item \emph{Parallel-jaw target.} The hand is reduced to a wrist pose and a
        single continuous gripper opening \eqref{eq:gripper}. Multi-fingered
        dexterous targets would require the finger-level representation
        rejected in Section~\ref{sec:lit-hand}.
  \item \emph{The gripper body is not the hand.} The target is the midpoint
        between thumb and index tips, but a parallel-jaw gripper interposes a
        body of its own between the flange and that point, and the whole
        assembly must stay above the workspace plane. For tasks performed
        against the table surface the demonstrated pose can therefore be
        kinematically unreachable even when the arm's reach is ample. The
        size of this gap requires evaluation on retained source episodes.
  \item \emph{Rigid objects.} The object-relative representation
        \eqref{eq:objrel} presupposes a rigid body with a well-defined pose.
        Deformable objects would require a flow-based representation
        \cite{novaflow}.
  \item \emph{Single-arm, quasi-static tasks.} The task class is restricted to
        single-arm tabletop manipulation at moderate speed, where the
        quasi-static assumption underlying purely kinematic retargeting holds.
  \item \emph{Fixed workspace.} The ChArUco anchor is rigidly attached to the
        workspace and the camera rig is not moved within a recording session.
  \item \emph{Kinematic-only transfer.} Forces and contact dynamics are not
        estimated. Tasks whose success depends on force regulation are outside
        the scope.
\end{enumerate}

\chapter{Implementation and Evaluation}
\label{ch:impl}

\section{Software and Experimental Setup}
\label{sec:arch}

The Python implementation separates capture, calibration, perception,
preparation, retargeting and export. FastAPI serves offline jobs and the
browser inspection interface; MuJoCo tools reside in a separate image.
Each stage writes a durable episode artifact. The \texttt{rec.npz} hand
extraction, \texttt{cln.npz} prepared motion,
\texttt{object\_models.npz} scene model and \texttt{plan.h5} robot plan
represent different stages and are not interchangeable evidence.
Physical-robot replay is a stub and is not evaluated.

The local apparatus has two Azure Kinect DK cameras on independent USB
controllers. The intended tabletop task is to transfer a blue cube from
a red support to a green target, using a simulated old-generation UR10
with a Robotiq 2F-85 gripper. No physical UR10 is available for this
study. The earlier local cube recordings and their derived artifacts
are not available for verification, so no numerical result from that
campaign is included in this manuscript. In particular, local
capture-synchronisation, stage-yield, object-model, IK, contact-replay
and learned-policy results are \emph{not} claimed. A new, preserved
local cohort would be required to evaluate that end-to-end chain.

\section{External HO-Cap Hand and Object Evaluation}
\label{sec:hocap-results}

The full external experiment covers 28 single-hand HO-Cap pick-and-place
sequences from nine subjects: 24 right-hand and four left-hand, with 29,706
frames from eight calibrated RealSense views. These are not the local Kinect
cube recordings. The source revision and all 114 numerical JSON files are preserved in
the paper's pre-retarget results snapshot with SHA-256 hashes. Capture
and object protocols are documented in the accompanying dated logs.
Both processing queues completed 28/28 scenes. This is a selected task cohort, not an evaluation
on the official HO-Cap held-out split \cite{hocap}.

\subsection{Hand-Joint Accuracy}

Every one of the 623,826 valid reference joint slots has a finite prepared
prediction. Matched-joint-weighted MPJPE is 10.79~mm and PCK@20/50~mm is
92.77/99.85\%. The per-sequence MPJPE range is 9.26--12.85~mm; the
per-sequence P95 range is 17.72--31.90~mm. No pooled P95 is inferred from
these 28 quantiles. The four left-hand scenes have weighted MPJPE
11.60~mm, versus 10.65~mm for 24 right-hand scenes; handedness and subject
are confounded here, so this is descriptive rather than a controlled
difference. The fused-observed stage, before final temporal smoothing,
matches 623,799/623,826 slots (99.996\%) with weighted MPJPE 11.16~mm.
The all-finite smoothed score includes 27 filled slots and should not be
mistaken for direct observation coverage. Table~\ref{tab:hocap-hand-summary}
summarizes the cohort; all 28 episode rows, including P95 and PCK, appear
in Table~\ref{tab:hocap-hand-detail}.

\begin{table}[htbp]
\caption{Eight-view HO-Cap 21-joint accuracy against published MANO fits}
\label{tab:hocap-hand-summary}
\centering\small
\begin{tabular}{lrrrr}
\toprule
Subset and stage & Scenes & Matched / GT & MPJPE (mm) & PCK@20 (\%) \\
\midrule
All, fused observed & 28 & 623,799/623,826 & 11.16 & 91.68 \\
All, smoothed all & 28 & 623,826/623,826 & 10.79 & 92.77 \\
Left, smoothed all & 4 & 86,667/86,667 & 11.60 & 87.56 \\
Right, smoothed all & 24 & 537,159/537,159 & 10.65 & 93.61 \\
\bottomrule
\end{tabular}
\end{table}

\subsection{Named Object Tracks and Outliers}

Each scene contains four annotated objects, giving 112 object--sequence
tracks across 15 object groups. The official HO-Cap object descriptors
\cite{hocap-catalog} are keyed by IDs in Table~\ref{tab:hocap-object-catalog};
for example, \texttt{G18\_2} is a projector remote, not an anonymous
object. All four instances in the \texttt{G18} scenes are the green tea
latte mix, projector remote, right shoe and left shoe. This identity
mapping comes from the authors' \texttt{hocap\_info.yaml}, rather than
appearance-based naming.

The tracker covers 117,916 of 118,824 valid reference object frames
(99.24\%); 97/112 tracks have complete coverage. Across the 112
\emph{track means}, the median and arithmetic mean relative translation
errors are 4.74 and 5.44~mm; the corresponding rotation values are
3.47 and 7.58 degrees. Frame-count-weighted averages of the track means
are 5.60~mm and 7.55 degrees. These are not pooled frame-level P95
values. The median of the 112 track P95s is 10.60~mm translation and
9.25 degrees rotation. First-frame alignment is fitted using the
reference pose, so the numbers measure subsequent relative drift rather
than absolute object-frame accuracy. The supplied instance masks
remove segmentation errors from this test; static intervals remain mixed
with manipulated intervals.

\begin{table}[htbp]
\caption{Selected HO-Cap object failures; mean and P95 are within-track errors}
\label{tab:hocap-object-outliers}
\centering\scriptsize
\begin{tabular}{llrrrrr}
\toprule
Sequence & Object and official descriptor & Cov. (\%) & T mean & T P95 & R mean & R P95 \\
\midrule
S7/162803 & \texttt{G18\_2} projector remote & 98.63 & 12.38 & 20.44 & 112.60 & 175.26 \\
S5/112303 & \texttt{G15\_2} dandelion tea & 100.00 & 17.24 & 21.81 & 109.65 & 128.43 \\
S2/202617 & \texttt{G09\_3} toothpaste box & 100.00 & 5.77 & 13.50 & 74.03 & 132.88 \\
S7/193506 & \texttt{G02\_4} sauce bottle & 99.65 & 2.82 & 10.67 & 56.57 & 169.89 \\
S2/163929 & \texttt{G19\_1} electric screwdriver & 48.80 & 5.75 & 15.41 & 10.55 & 50.95 \\
S2/201942 & \texttt{G06\_1} DVD & 100.00 & 19.39 & 25.24 & 1.90 & 5.63 \\
\bottomrule
\end{tabular}
\end{table}

The worst rotation case is the projector remote in
\texttt{subject\_7/162803}: mean 112.60 degrees and P95 175.26 degrees
despite 98.63\% tracked coverage and 12.38~mm mean translation error.
In another \texttt{G18} scene (\texttt{subject\_1/165502}), the same
object ID has 6.97-degree mean rotation error, and the other three
objects in the failing scene have 1.10--1.34-degree means. Thus the
anomaly is sequence/object-specific rather than a global frame rotation.
A visual review showed plausible late object placement under an alternate
orientation, but did not establish a unique root cause; the outlier remains
in all aggregates. The dandelion tea (109.65 degrees) and toothpaste box
(74.03 degrees in one scene, versus 13.82 and 6.89 degrees in two others)
show that this is not a single isolated failure. A low translation error
also does not certify orientation, as the sauce bottle demonstrates.
Conversely, the DVD has the highest mean translation error, 19.39~mm,
with only 1.90-degree mean rotation error. The electric screwdriver is a
coverage failure at 48.80\%, so its conditional error on visible frames
must not stand in for whole-sequence accuracy. Tables
\ref{tab:hocap-object-detail} and \ref{tab:hocap-object-catalog} retain
all 112 object scores and official identities for audit.

\subsection{Historical Two-Sequence Camera-Pair Pilot}

Before the full eight-view run, two subject-1 sequences were tested for
120 frames each with a pair selected after screening all 28 pairs on those
same frames. Table~\ref{tab:hocap-hand} is retained as a historical
\emph{in-sample} development comparison; its original local artifacts
were lost. The full-sequence eight-view results above supersede it as the
reproducible hand/object reference cohort, but do not turn the selected
pair into a held-out baseline.

\begin{table}[htbp]
\caption{Historical HO-Cap Hand-Joint Pilot: Selected Pair Versus Eight Views}
\label{tab:hocap-hand}
\centering\small
\begin{tabular}{llrrrr}
\toprule
Sequence & Views & Coverage & MPJPE (mm) & P95 (mm) & PCK@20 (\%) \\
\midrule
\texttt{170650} & selected pair & 2,520/2,520 & 11.33 & 22.71 & 87.70 \\
\texttt{170650} & all eight & 2,520/2,520 & 11.34 & 22.17 & 87.50 \\
\texttt{170105} & selected pair & 2,520/2,520 & 11.67 & 22.80 & 89.80 \\
\texttt{170105} & all eight & 2,520/2,520 & 10.50 & 20.15 & 94.88 \\
\bottomrule
\end{tabular}
\end{table}

The selected pair also failed one colour--depth consistency check. No
HO-Cap score establishes the accuracy of the local Kinect calibration or
physical success of the proposed robot motion.


\section{Downstream Evidence Boundary}
\label{sec:downstream-boundary}

The archived HO-Cap experiment evaluates hand reconstruction and
first-frame-aligned object tracking on a separate RealSense capture
system. It does not establish the accuracy of the local Kinect rig or
produce a validated robot demonstration. Retarget plans, simulator
contacts, task success and policy learning require their own retained
inputs, protocols and numerical artifacts before they can be reported
as results. This chapter therefore ends at the external perception
evaluation rather than extrapolating an end-to-end success rate.

\chapter{Analysis and Discussion}
\label{ch:disc}

\section{What the Results Establish}
\label{sec:interpretation}

The retained eight-view HO-Cap evaluation covers 28 selected
pick-and-place sequences from nine subjects. Prepared 21-joint hand
predictions match all 623,826 valid reference joint slots, with
10.79~mm matched-joint-weighted MPJPE against the dataset's published
semi-automatic MANO fits. This is an external-system perception
measurement, not evidence that the local Azure Kinect rig has the
same accuracy. The object evaluation covers 112 named tracks, but it
uses published instance masks and a reference-assisted first-frame
SE(3) alignment. Its 5.44~mm translation and 7.58-degree rotation
macro-means describe relative tracking drift conditional on those
inputs, not absolute object localization or autonomous segmentation.

\section{Error Sources and Reference Frames}
\label{sec:errbudget}

The hand and object numbers cannot be added into an end-to-end geometric
error budget. The hand score is an absolute comparison in the HO-Cap
world frame; the object score removes a per-object initial gauge
difference. Both refer to semi-automatic annotations rather than an
independent metrology system. The object tracker also receives masks
generated outside ViKi, and static intervals are included alongside
manipulation. Thus a small average object drift need not imply a
sufficiently accurate grasp pose during contact.

The outliers make this distinction operational. The projector remote
\texttt{G18\_2} has 112.60-degree mean rotation error in one sequence
at 98.63\% coverage, but only 6.97 degrees in another. The dandelion
tea and toothpaste box have additional large orientation errors.
The sauce bottle combines 2.82~mm mean translation error with
56.57-degree mean orientation error. These tracks remain in the
aggregate rather than being removed as failures. Symmetry-aware
orientation scoring was not used; the visual review of the remote
did not uniquely diagnose the failure. Conversely, the electric
screwdriver is tracked for only 48.80\% of valid reference frames,
so its conditional error does not describe the missing half.

For the local two-Kinect system, camera geometry, RGB--depth alignment,
exposure timing and table-height registration remain unquantified in
this manuscript. The former local numerical summaries cannot be
rechecked because their source recordings and derived artifacts are
unavailable. Code paths and visually plausible clouds are not
substitutes for a preserved evaluation protocol.

\section{Statistical Validity and Attribution}
\label{sec:validity}

The HO-Cap selection is a single-hand pick-and-place subset, not the
official held-out benchmark split. Subjects and hands are unevenly
represented: 24 right-hand and four left-hand sequences. Joint slots
within the same frames and object frames within the same sequence
are correlated, so their large counts do not imply that many
independent trials. The two-sequence camera-pair pilot was selected
on the same development frames and cannot be interpreted as held-out
evidence. The main eight-view aggregate is descriptive for this
selected cohort; no causal claim about the benefit of camera count,
training data or robot embodiment follows from it.

No surviving local cube dataset supports the previous downstream
figures. In particular, there is no verifiable local stage yield,
simulated stacking rate, policy comparison or contact-safety result.
The current manuscript therefore does not claim an end-to-end
camera-to-simulation proof of concept or Real-to-Sim-to-Real transfer.

\section{Scope and Implications}
\label{sec:implications}

The supported contribution is a reproducible external perception
assessment, with explicit object identities, coverage, outliers and
metric limitations, alongside an implemented offline software
pipeline. The intended cube-transfer experiment remains a research
question rather than a completed result. A new preserved local
cohort should include measurement-only controls: known geometry at
table and carry heights, a marked asymmetric object for orientation,
and disjoint human trials. Retained simulator rollouts and matched
no-demonstration baselines would then be needed before reporting
physical task success or learning benefit. No physical UR10
evaluation is part of the current evidence.

\chapter{Conclusion}
\label{ch:concl}

\section{Summary}
\label{sec:summary}

ViKi converts passive two-view RGB-D recordings of a human cube transfer
into separate hand, object and proposed UR10 trajectory artifacts. A
geometric plan is not itself a robot demonstration; dynamic replay with
free objects supplies the stricter simulation test. This thesis reports
the measured chain on 78 recordings and a local layout-conditioned
learning experiment rather than promising a later unmeasured result.

\section{Answers to the Research Questions}
\label{sec:answers}

\textbf{RQ1: pipeline yield.} Seventy-three of 78 recordings pass the
articulated hand gate; 44 of those pass the original object-quality gate
and 70 pass after independent scratch track recovery. Of 26 newly
recovered scenes, 22 pass geometric retargeting. Seventeen of those 22
also produce released stable stacks in the provisional MuJoCo replay.
The branch-specific denominators and the red-support contact in all 24
replayed scratch plans preclude one unqualified success percentage.
The external HO-Cap pilot measured 10.50--11.67~mm smoothed hand-joint MPJPE
against semi-automatic annotations on two short sequences using either the
selected pair or all eight views; it does not ground-truth the local cube
recordings or their object and robot trajectories. Its source episodes and
the local cohort were lost after the pilot and must be reacquired for a
reproducible final evaluation.

\textbf{RQ2: local demonstration generation.} One recovered human
motion generated 46 successful simulator variants in a 48-layout
$\pm20$~mm sweep and 27 in an independent 48-layout $\pm50$~mm sweep.
All derive from that one motion and touch the red support. This is
local, physics-filtered candidate generation, not arbitrary scene or
robot-configuration transfer.

\textbf{RQ3: learned correction.} A phase- and initial-layout-conditioned
model trained on 73 generated variants outperformed unchanged source
commands on the same 14 held-out teacher-selected layouts (14/14 versus
4/14) and 30 separately sampled layouts (30/30 versus 0/30). The earlier
plain state BC failed in all 17 closed-loop trials. The successful model
retains the source trajectory and observes privileged initial cube
positions; it is not RL or a vision-based closed-loop policy.

\section{Claim Boundary and Next Validation}
\label{sec:contrib-final}

The evidence supports a camera-to-simulation proof of concept for one
human source trajectory, local cube-layout changes, one old-UR10/Robotiq
model and one provisional contact configuration. It does not demonstrate
collision-free grasping, independent 3-D orientation accuracy,
source-disjoint generalization, visual-policy performance, benefit over
a separately trained no-data learner or real-robot transfer. Those
claims require additional tests rather than a larger number of
synthetic variants of the same source motion.

A small measurement-only set with known geometry, marked orientation
and independently held-out human trials is the practical next step
for the thesis evidence. Support-object clearance and contact-model
sensitivity need correction before treating a simulated success as a
robot-ready reference. A physical UR10 evaluation would additionally
require camera-to-robot registration and hardware safety validation;
it is not part of the present result.

\section{Concluding Remarks}
\label{sec:remarks}

The result is useful precisely because the stages are not conflated:
recording completeness, reconstruction coverage, model acceptance,
geometric IK, dynamic stack success and learning on new layouts are
reported separately. The local learning result is positive; the
red-support collisions and failed plain BC are equally important
boundaries on what ViKi has actually established.

\chapter*{References}
\addcontentsline{toc}{chapter}{References}

\begin{thebibliography}{99}
\setlength{\itemsep}{2pt}

\bibitem{mimicgen}
A. Mandlekar, S. Nasiriany, B. Wen, I. Akinola, Y. Narang, L. Fan, Y. Zhu,
and D. Fox, ``MimicGen: A Data Generation System for Scalable Robot Learning
using Human Demonstrations,'' in \emph{Proc. 7th Conf. on Robot Learning},
PMLR 229, pp. 1820--1864, 2023.
[Online]. Available: \url{https://proceedings.mlr.press/v229/mandlekar23a.html}

\bibitem{anyteleop}
Y. Qin, W. Yang, B. Huang, K. Van Wyk, H. Su, X. Wang, Y.-W. Chao, and D. Fox,
``AnyTeleop: A general vision-based dexterous robot arm-hand teleoperation
system,'' in \emph{Proc. Robotics: Science and Systems (RSS)}, 2023.
[Online]. Available: \url{https://arxiv.org/abs/2307.04577}

\bibitem{opentv}
X. Cheng, J. Li, S. Yang, G. Yang, and X. Wang, ``Open-TeleVision:
Teleoperation with immersive active visual feedback,'' in \emph{Proc. Conf. on
Robot Learning (CoRL)}, 2024. [Online]. Available:
\url{https://arxiv.org/abs/2407.01512}

\bibitem{bunny}
R. Ding, Y. Qin, J. Zhu, C. Jia, S. Yang, R. Yang, X. Qi, and X. Wang,
``Bunny-VisionPro: Real-time bimanual dexterous teleoperation for imitation
learning,'' in \emph{Proc. IEEE/RSJ Int. Conf. on Intelligent Robots and
Systems (IROS)}, 2025. [Online]. Available:
\url{https://arxiv.org/abs/2407.03162}

\bibitem{ace}
S. Yang, M. Liu, Y. Qin, R. Ding, J. Li, X. Cheng, R. Yang, S. Yi, and X. Wang,
``ACE: A cross-platform visual-exoskeletons system for low-cost dexterous
teleoperation,'' in \emph{Proc. Conf. on Robot Learning (CoRL)}, 2024.
[Online]. Available: \url{https://arxiv.org/abs/2408.11805}

\bibitem{dexflow}
\ph{[Authors]}, ``DexFlow: Dexterous retargeting with $C^2$ continuity via
Hessian regularisation,'' \emph{arXiv:2505.01083}, 2025. \ph{[Verify author
list and venue.]}

\bibitem{dexteleop0}
\ph{[Authors]}, ``\ph{[Title]},'' 2025. \ph{[Robust Huber-loss retargeting
formulation cited in Section~\ref{sec:lit-retarget}; confirm identifier before
submission.]}

\bibitem{keyobj}
\ph{[C.]} Xin, \ph{[Y.]} Yu, \ph{[X.]} Jiang, \ph{[Y.]} Zhang, and \ph{[H.]}
Li, ``Analyzing key objectives in human-to-robot retargeting for dexterous
manipulation,'' \emph{arXiv:2506.09384}, 2025.

\bibitem{pink}
S. Caron \emph{et al.}, ``Pink: Python inverse kinematics based on Pinocchio,''
software library. [Online]. Available:
\url{https://github.com/stephane-caron/pink}

\bibitem{orion}
Y. Zhu, \ph{[et al.]}, ``Vision-based manipulation from single human video with
open-world object graphs,'' \emph{arXiv:2405.20321}, 2024; \emph{Autonomous
Robots}, 2026. \ph{[Confirm final volume and pages.]}

\bibitem{okami}
J. Li, Y. Zhu, Y. Xie, Z. Jiang, M. Seo, G. Pavlakos, and Y. Zhu, ``OKAMI:
Teaching humanoid robots manipulation skills through single video imitation,''
in \emph{Proc. Conf. on Robot Learning (CoRL)}, \emph{PMLR}, vol.~270,
pp.~299--317, 2024.

\bibitem{dexmv}
Y. Qin, Y.-H. Wu, S. Liu, H. Jiang, R. Yang, Y. Fu, and X. Wang, ``DexMV:
Imitation learning for dexterous manipulation from human videos,'' in
\emph{Proc. European Conf. on Computer Vision (ECCV)}, 2022.

\bibitem{spot}
\ph{[Authors]}, ``SPOT: SE(3) pose trajectory diffusion for object-centric
manipulation,'' in \emph{Proc. IEEE Int. Conf. on Robotics and Automation
(ICRA)}, 2025. [Online]. Available: \url{https://arxiv.org/abs/2411.00965}

\bibitem{novaflow}
\ph{[Authors]}, ``NovaFlow: Zero-shot manipulation via actionable flow from
generated videos,'' \emph{arXiv:2510.08568}, 2025.

\bibitem{demoscore}
A. Chen, A. Lessing, Y. Liu, and C. Finn, ``Curating demonstrations using
online experience,'' \emph{arXiv:2503.03707}, 2025.

\bibitem{openvlaoft}
M. J. Kim, C. Finn, and P. Liang, ``Fine-tuning vision-language-action models:
Optimizing speed and success,'' \emph{arXiv:2502.19645}, 2025.

\bibitem{egomimic}
S. Kareer, D. Patel, R. Punamiya, P. Mathur, S. Cheng, C. Wang, J. Hoffman, and
D. Xu, ``EgoMimic: Scaling imitation learning via egocentric video,'' in
\emph{Proc. Conf. on Robot Learning (CoRL)}, 2024. [Online]. Available:
\url{https://arxiv.org/abs/2410.24221}

\bibitem{dexwild}
T. Tao, M. K. Srirama, J. J. Liu, K. Shaw, and D. Pathak, ``DexWild: Dexterous
human interactions for in-the-wild robot policies,'' in \emph{Proc. Robotics:
Science and Systems (RSS)}, 2025. [Online]. Available:
\url{https://arxiv.org/abs/2505.07813}

\bibitem{toolif}
\ph{[Authors]}, ``Tool-as-Interface: Learning robot policies from observing
human tool use,'' \emph{arXiv:2504.04612}, 2025.

\bibitem{howtransfer}
\ph{[Authors]}, ``Hand-centric human-to-robot trajectory transfer from video
demonstrations via open-world contact localization,'' \emph{arXiv:2606.10743},
2026. \ph{[Preprint; peer-review status unconfirmed.]}

\bibitem{viroil}
\ph{[Authors]}, ``Video-to-robot imitation learning for bimanual dexterous
manipulation from human video,'' in \emph{Proc. ICEILM}, 2025,
doi:10.1145/3788731.3788736. \ph{[Verify author list from the published PDF.]}

\bibitem{robopaint}
\ph{[Authors]}, ``RoboPaint: From human demonstration to any robot and any
view,'' \emph{arXiv:2602.05325}, 2026. \ph{[Preprint; status unconfirmed.]}

\bibitem{shadowing}
\ph{[Authors]}, ``Vision-based hand shadowing,'' \emph{IEEE Access}, 2026;
\emph{arXiv:2603.11383}. \ph{[Preprint; status unconfirmed.]}

\bibitem{dexcanvas}
\ph{[Authors]}, ``DexCanvas: Bridging human demonstrations and robot learning
for dexterous manipulation,'' \emph{arXiv:2510.15786}, 2025.

\bibitem{romeo}
L. Romeo, R. Marani, A. G. Perri, and T. D'Orazio, ``Microsoft Azure Kinect
calibration for three-dimensional dense point clouds and reliable skeletons,''
\emph{Sensors}, vol.~22, no.~13, art.~4986, 2022.

\bibitem{darwish}
W. Darwish, Q. Bols\'{e}e, and A. Munteanu, ``Robust calibration of a
multi-view Azure Kinect scanner based on spatial consistency,'' in
\emph{Proc. IEEE Int. Conf. on 3D Immersion (IC3D)}, 2020.

\bibitem{kimcal}
J. Kim \emph{et al.}, ``Robust extrinsic calibration of multiple RGB-D cameras
with body tracking and feature matching,'' \emph{Sensors}, vol.~21, no.~3,
art.~1013, 2021.

\bibitem{bulee}
S. Bu and S. Lee, ``Easy to calibrate: Marker-less calibration of multiview
Azure Kinect,'' \emph{Computer Modeling in Engineering \& Sciences},
vol.~136, no.~3, pp.~3083--3096, 2023.

\bibitem{tolgyessy}
M. T\"{o}lgyessy, M. Dekan, \v{L}. Chovanec, and P. Hubinsk\'{y}, ``Evaluation
of the Azure Kinect and its comparison to Kinect v1 and Kinect v2,''
\emph{Sensors}, vol.~21, no.~2, art.~413, 2021.

\bibitem{kurillo}
G. Kurillo, E. Hemingway, M.-L. Cheng, and L. Cheng, ``Evaluating the accuracy
of the Azure Kinect and Kinect v2,'' \emph{Sensors}, vol.~22, no.~7, art.~2469,
2022.

\bibitem{hocap}
J. Wang, Q. Zhang, Y.-W. Chao, B. Wen, X. Guo, and Y. Xiang,
HO-Cap: A Capture System and Dataset for 3D Reconstruction and Pose
Tracking of Hand-Object Interaction, in \emph{Advances in Neural Information
Processing Systems, Datasets and Benchmarks Track}, 2025.
[Online]. Available: \url{https://irvlutd.github.io/HOCap/}
\bibitem{hocap-catalog}
IRVLUTD, {HO-Cap Toolkit object descriptors},
\texttt{config/hocap\_info.yaml}.
[Online]. Available:
\url{https://github.com/IRVLUTD/HO-Cap/blob/main/config/hocap_info.yaml}.
Accessed: 14 October 2026.

\bibitem{mskinect}
Microsoft, ``Azure Kinect DK hardware specification'' and ``Synchronize
multiple Azure Kinect DK devices,'' Microsoft Learn documentation.
\ph{[Add access date.]}

\bibitem{intel}
Intel Corporation, ``Intel RealSense D400 series product family datasheet'' and
multi-camera configuration guidance. \ph{[Add revision and access date.]}

\bibitem{ridgerun}
RidgeRun, ``Multi-camera interface configurations and synchronization,''
technical note. \ph{[Non-peer-reviewed vendor documentation; retain only for
engineering practice, not for scientific claims.]}

\bibitem{alihandeye}
I. Ali \emph{et al.}, ``Accuracy evaluation of hand-eye calibration techniques
for vision-guided robots,'' \emph{PLOS ONE}, 2022. \ph{[Add volume and article
number.]}

\bibitem{fusionpatent}
``Image-based geometric fusion of multiple depth images using ray casting,''
U.S. Patent 10\,818\,071, 2020.

\bibitem{zeng}
J. Zeng, Y. Tong, Y. Huang, Q. Yan, W. Sun, J. Chen, and Y. Wang, ``Deep
surface normal estimation with hierarchical RGB-D fusion,'' in \emph{Proc. IEEE
Conf. on Computer Vision and Pattern Recognition (CVPR)}, 2019. [Online].
Available: \url{https://arxiv.org/abs/1904.03405}

\bibitem{hrbf}
Y. Wang \emph{et al.}, ``HRBF-Fusion: Accurate 3D reconstruction from RGB-D
data using on-the-fly implicits,'' \emph{ACM Transactions on Graphics}.
\ph{[Add volume, issue and year.]}

\end{thebibliography}

\vspace{4mm}
\noindent\ph{\small Note: entries marked in this colour require completion or
verification before submission. Preprint entries dated 2026 could not be
confirmed as peer-reviewed at the time of writing and are cited as preprints.}

\appendix
\chapter{Notation}
\label{app:notation}

\begin{table}[htbp]
\caption{Symbols Used Throughout the Thesis}
\label{tab:notation}
\centering
\small
\begin{tabular}{ll}
\toprule
Symbol & Meaning \\
\midrule
$W$ & workspace frame, defined by the ChArUco anchor \\
$C_k$ & frame of camera $k$ \\
$R$ & manipulator base frame \\
$E$ & manipulator end-effector frame \\
$H$ & human wrist frame \\
$O$ & manipulated object frame \\
$\bT^{A}_{B}$ & rigid transform from frame $B$ to frame $A$, element of $\SE$ \\
$\bq_t \in \R^{n}$ & manipulator configuration at timestep $t$ \\
$f(\bq_t)$ & forward kinematics of the end-effector \\
$\be_t$ & task-space residual, \eqref{eq:residual} \\
$\omega_t$ & per-frame perception confidence weight, \eqref{eq:conf} \\
$q_{i,t}$ & geometric quality of a triangulated landmark, \eqref{eq:quality} \\
$w^{k}_{i,t}$ & per-view weight, single-view fallback, \eqref{eq:fusion} \\
$v^{k}_{i,t}$ & detector visibility score \\
$c^{k}_{i,t}$ & sensor depth confidence \\
$\rho_\delta$ & Huber loss with transition parameter $\delta$ \\
$\lambda_v,\lambda_a,\lambda_r$ & velocity, acceleration and posture weights \\
$g_t$ & continuous gripper opening, \eqref{eq:gripper} \\
$\tau$ & host-side monotonic timestamp, \eqref{eq:hostclock} \\
$T$ & number of timesteps in an episode \\
\bottomrule
\end{tabular}
\end{table}

\chapter{Full HO-Cap Perception Results}
\label{app:hocap}
Sequence labels use S<subject>/<HHMMSS> from the original HO-Cap identifier.
The complete unrounded JSON and checksums are preserved in
\texttt{paper/results/hocap/pre\_retarget/}. Tables in this appendix can
be regenerated with \texttt{python3 scripts/build\_hocap\_paper\_tables.py};
the script refuses modified source hashes or an incomplete cohort.
Object descriptors follow the official HO-Cap toolkit's
\texttt{config/hocap\_info.yaml} \cite{hocap-catalog}. Hand P95 is per sequence;
object P95 is per object--sequence track, never pooled across rows.

\begin{longtable}{lrrrrrr}
\caption{Eight-view HO-Cap hand results by sequence; PCK values are percentages.}\label{tab:hocap-hand-detail}\\
\toprule Sequence & Frames & Matched & MPJPE & P95 & PCK@20 & PCK@50 \\
\midrule\endfirsthead
\toprule Sequence & Frames & Matched & MPJPE & P95 & PCK@20 & PCK@50 \\
\midrule\endhead
S1/165502 & 676 & 14196 & 12.85 & 25.48 & 82.83 & 99.87 \\
S1/165807 & 798 & 16758 & 12.10 & 24.18 & 84.35 & 99.79 \\
S1/170105 & 831 & 17451 & 11.61 & 21.46 & 92.71 & 99.91 \\
S1/170650 & 741 & 15561 & 12.55 & 23.78 & 86.36 & 100.00 \\
S1/170959 & 874 & 18354 & 11.26 & 21.67 & 92.45 & 99.96 \\
S2/201316 & 828 & 17388 & 10.77 & 21.23 & 93.48 & 99.95 \\
S2/201942 & 1561 & 32781 & 11.86 & 23.79 & 86.87 & 99.95 \\
S2/202617 & 1271 & 26691 & 11.06 & 22.30 & 91.43 & 99.99 \\
S2/163929 & 1129 & 23709 & 10.48 & 20.43 & 94.11 & 99.99 \\
S3/154531 & 1686 & 35406 & 9.26 & 18.09 & 96.95 & 99.98 \\
S3/155008 & 1046 & 21966 & 9.45 & 17.72 & 97.89 & 100.00 \\
S3/161209 & 1071 & 22491 & 9.71 & 18.75 & 96.10 & 99.92 \\
S3/161937 & 1269 & 26649 & 10.33 & 19.48 & 95.74 & 99.79 \\
S3/162327 & 878 & 18438 & 9.82 & 19.30 & 95.79 & 99.98 \\
S3/162756 & 946 & 19866 & 9.56 & 17.79 & 98.23 & 100.00 \\
S4/163223 & 939 & 19719 & 12.54 & 30.34 & 86.51 & 99.24 \\
S4/164812 & 974 & 20454 & 10.57 & 20.61 & 94.22 & 99.99 \\
S5/112303 & 702 & 14742 & 10.95 & 26.74 & 91.05 & 99.54 \\
S6/110646 & 1292 & 27132 & 10.02 & 19.26 & 95.51 & 99.98 \\
S6/111118 & 1272 & 26712 & 10.54 & 21.84 & 93.52 & 99.87 \\
S6/112546 & 1311 & 27531 & 10.06 & 21.66 & 93.83 & 99.69 \\
S7/193506 & 1136 & 23856 & 12.40 & 31.90 & 87.20 & 98.66 \\
S7/162803 & 1092 & 22932 & 10.10 & 21.11 & 93.81 & 99.86 \\
S8/180651 & 985 & 20685 & 11.88 & 24.84 & 89.59 & 99.88 \\
S9/123403 & 992 & 20832 & 10.28 & 19.54 & 95.56 & 99.99 \\
S9/123725 & 1134 & 23814 & 10.95 & 20.10 & 94.81 & 99.99 \\
S9/124057 & 1136 & 23856 & 10.88 & 21.39 & 93.45 & 99.84 \\
S9/125315 & 1136 & 23856 & 10.74 & 20.45 & 94.53 & 99.97 \\
\bottomrule
\end{longtable}

\begin{longtable}{llrrrrr}
\caption{Per-object relative-pose tracking after one first-frame SE(3) alignment. T and R are translation in mm and geodesic rotation in degrees.}\label{tab:hocap-object-detail}\\
\toprule Sequence & Object & Cov. (\%) & T mean & T P95 & R mean & R P95 \\
\midrule\endfirsthead
\toprule Sequence & Object & Cov. (\%) & T mean & T P95 & R mean & R P95 \\
\midrule\endhead
S1/165502 & G18\_1 & 100.00 & 2.68 & 3.75 & 1.32 & 1.97 \\
S1/165502 & G18\_2 & 100.00 & 7.70 & 24.15 & 6.97 & 27.41 \\
S1/165502 & G18\_3 & 100.00 & 1.49 & 6.97 & 1.32 & 7.71 \\
S1/165502 & G18\_4 & 100.00 & 1.66 & 7.76 & 2.20 & 5.86 \\
S1/165807 & G15\_1 & 100.00 & 2.41 & 10.53 & 1.58 & 3.85 \\
S1/165807 & G15\_2 & 100.00 & 2.00 & 2.76 & 3.02 & 3.34 \\
S1/165807 & G15\_3 & 100.00 & 2.17 & 10.08 & 1.24 & 7.77 \\
S1/165807 & G15\_4 & 100.00 & 6.02 & 12.56 & 10.65 & 18.60 \\
S1/170105 & G10\_1 & 100.00 & 0.67 & 3.56 & 1.55 & 10.44 \\
S1/170105 & G10\_2 & 100.00 & 5.57 & 6.37 & 11.96 & 20.04 \\
S1/170105 & G10\_3 & 100.00 & 10.13 & 23.41 & 3.54 & 9.66 \\
S1/170105 & G10\_4 & 100.00 & 1.23 & 4.46 & 1.06 & 3.36 \\
S1/170650 & G04\_1 & 100.00 & 1.41 & 4.74 & 1.64 & 3.84 \\
S1/170650 & G04\_2 & 100.00 & 2.21 & 3.05 & 1.08 & 2.63 \\
S1/170650 & G04\_3 & 100.00 & 8.89 & 20.10 & 6.02 & 14.70 \\
S1/170650 & G04\_4 & 100.00 & 0.72 & 2.98 & 0.32 & 2.15 \\
S1/170959 & G19\_1 & 100.00 & 2.41 & 8.12 & 4.74 & 15.42 \\
S1/170959 & G19\_2 & 100.00 & 12.85 & 46.05 & 3.41 & 8.96 \\
S1/170959 & G19\_3 & 100.00 & 2.95 & 10.80 & 5.50 & 19.73 \\
S1/170959 & G19\_4 & 100.00 & 14.36 & 28.21 & 13.02 & 65.53 \\
S2/201316 & G05\_1 & 100.00 & 4.73 & 12.06 & 3.30 & 9.27 \\
S2/201316 & G05\_2 & 100.00 & 0.94 & 4.74 & 0.74 & 4.64 \\
S2/201316 & G05\_3 & 100.00 & 9.40 & 28.11 & 7.81 & 17.53 \\
S2/201316 & G05\_4 & 100.00 & 6.03 & 12.08 & 4.18 & 12.11 \\
S2/201942 & G06\_1 & 100.00 & 19.39 & 25.24 & 1.90 & 5.63 \\
S2/201942 & G06\_2 & 100.00 & 0.55 & 4.18 & 0.67 & 7.04 \\
S2/201942 & G06\_3 & 100.00 & 6.42 & 10.68 & 7.26 & 17.51 \\
S2/201942 & G06\_4 & 99.68 & 7.40 & 28.59 & 13.03 & 40.35 \\
S2/202617 & G09\_1 & 100.00 & 11.27 & 24.44 & 3.91 & 18.56 \\
S2/202617 & G09\_2 & 100.00 & 3.93 & 8.17 & 3.14 & 6.62 \\
S2/202617 & G09\_3 & 100.00 & 5.77 & 13.50 & 74.03 & 132.88 \\
S2/202617 & G09\_4 & 99.84 & 13.28 & 31.71 & 5.18 & 12.13 \\
S2/163929 & G19\_1 & 48.80 & 5.75 & 15.41 & 10.55 & 50.95 \\
S2/163929 & G19\_2 & 100.00 & 9.67 & 25.64 & 1.94 & 5.60 \\
S2/163929 & G19\_3 & 100.00 & 2.13 & 7.22 & 6.38 & 8.89 \\
S2/163929 & G19\_4 & 100.00 & 10.74 & 34.86 & 2.20 & 6.96 \\
S3/154531 & G09\_1 & 100.00 & 5.93 & 13.05 & 2.77 & 9.75 \\
S3/154531 & G09\_2 & 100.00 & 6.10 & 13.35 & 6.79 & 28.73 \\
S3/154531 & G09\_3 & 100.00 & 9.03 & 14.95 & 13.82 & 142.95 \\
S3/154531 & G09\_4 & 100.00 & 7.41 & 17.99 & 4.77 & 7.02 \\
S3/155008 & G10\_1 & 100.00 & 1.59 & 4.16 & 1.93 & 7.28 \\
S3/155008 & G10\_2 & 100.00 & 5.55 & 12.66 & 3.94 & 17.99 \\
S3/155008 & G10\_3 & 100.00 & 5.67 & 7.58 & 1.70 & 6.88 \\
S3/155008 & G10\_4 & 100.00 & 4.51 & 8.15 & 2.31 & 9.23 \\
S3/161209 & G16\_1 & 100.00 & 1.91 & 4.31 & 5.78 & 6.44 \\
S3/161209 & G16\_2 & 100.00 & 4.74 & 11.53 & 13.93 & 90.87 \\
S3/161209 & G16\_3 & 100.00 & 2.86 & 7.87 & 1.92 & 10.59 \\
S3/161209 & G16\_4 & 100.00 & 9.37 & 20.60 & 9.55 & 41.50 \\
S3/161937 & G19\_1 & 100.00 & 3.11 & 16.36 & 7.25 & 51.05 \\
S3/161937 & G19\_2 & 100.00 & 5.98 & 14.12 & 1.52 & 5.72 \\
S3/161937 & G19\_3 & 100.00 & 3.50 & 7.80 & 7.81 & 18.02 \\
S3/161937 & G19\_4 & 100.00 & 18.47 & 45.93 & 4.65 & 12.46 \\
S3/162327 & G20\_1 & 100.00 & 0.99 & 2.36 & 0.93 & 2.31 \\
S3/162327 & G20\_2 & 100.00 & 4.66 & 10.04 & 3.04 & 7.04 \\
S3/162327 & G20\_3 & 100.00 & 1.93 & 9.06 & 0.62 & 2.82 \\
S3/162327 & G20\_4 & 95.22 & 6.53 & 14.24 & 6.58 & 21.72 \\
S3/162756 & G21\_1 & 100.00 & 2.51 & 4.55 & 2.18 & 3.26 \\
S3/162756 & G21\_2 & 100.00 & 3.91 & 14.02 & 8.22 & 52.38 \\
S3/162756 & G21\_3 & 100.00 & 10.18 & 22.65 & 2.89 & 6.76 \\
S3/162756 & G21\_4 & 99.89 & 4.70 & 10.28 & 4.46 & 29.13 \\
S4/163223 & G16\_1 & 100.00 & 5.17 & 23.66 & 5.12 & 25.23 \\
S4/163223 & G16\_2 & 100.00 & 3.78 & 9.71 & 4.00 & 9.65 \\
S4/163223 & G16\_3 & 100.00 & 3.58 & 5.05 & 10.68 & 15.81 \\
S4/163223 & G16\_4 & 93.93 & 9.75 & 28.59 & 9.55 & 35.11 \\
S4/164812 & G21\_1 & 100.00 & 1.72 & 3.03 & 1.23 & 3.54 \\
S4/164812 & G21\_2 & 100.00 & 2.86 & 10.76 & 6.74 & 16.99 \\
S4/164812 & G21\_3 & 100.00 & 15.06 & 32.43 & 3.60 & 10.91 \\
S4/164812 & G21\_4 & 100.00 & 7.72 & 13.79 & 5.29 & 19.97 \\
S5/112303 & G15\_1 & 100.00 & 1.12 & 5.40 & 1.91 & 13.57 \\
S5/112303 & G15\_2 & 100.00 & 17.24 & 21.81 & 109.65 & 128.43 \\
S5/112303 & G15\_3 & 100.00 & 5.28 & 8.30 & 8.28 & 12.34 \\
S5/112303 & G15\_4 & 100.00 & 4.63 & 11.75 & 8.16 & 27.42 \\
S6/110646 & G05\_1 & 100.00 & 2.76 & 6.17 & 1.83 & 5.31 \\
S6/110646 & G05\_2 & 100.00 & 1.77 & 3.99 & 0.69 & 1.72 \\
S6/110646 & G05\_3 & 91.33 & 5.27 & 21.45 & 4.69 & 28.75 \\
S6/110646 & G05\_4 & 98.22 & 7.22 & 17.21 & 9.47 & 24.79 \\
S6/111118 & G06\_1 & 99.37 & 10.02 & 18.20 & 2.08 & 7.79 \\
S6/111118 & G06\_2 & 100.00 & 1.49 & 2.64 & 5.78 & 8.36 \\
S6/111118 & G06\_3 & 100.00 & 1.33 & 3.50 & 2.61 & 6.06 \\
S6/111118 & G06\_4 & 96.31 & 8.01 & 13.28 & 41.23 & 65.31 \\
S6/112546 & G22\_1 & 99.31 & 6.76 & 15.41 & 8.75 & 82.09 \\
S6/112546 & G22\_2 & 100.00 & 1.77 & 4.50 & 1.49 & 2.46 \\
S6/112546 & G22\_3 & 100.00 & 0.60 & 2.39 & 0.95 & 2.31 \\
S6/112546 & G22\_4 & 100.00 & 10.98 & 30.93 & 2.42 & 7.98 \\
S7/193506 & G02\_1 & 100.00 & 7.51 & 8.67 & 2.30 & 5.40 \\
S7/193506 & G02\_2 & 99.65 & 4.85 & 7.52 & 4.64 & 11.53 \\
S7/193506 & G02\_3 & 100.00 & 1.10 & 5.58 & 1.12 & 6.37 \\
S7/193506 & G02\_4 & 99.65 & 2.82 & 10.67 & 56.57 & 169.89 \\
S7/162803 & G18\_1 & 100.00 & 5.49 & 7.02 & 1.33 & 4.87 \\
S7/162803 & G18\_2 & 98.63 & 12.38 & 20.44 & 112.60 & 175.26 \\
S7/162803 & G18\_3 & 100.00 & 0.98 & 6.86 & 1.10 & 6.94 \\
S7/162803 & G18\_4 & 100.00 & 1.64 & 4.22 & 1.34 & 4.97 \\
S8/180651 & G06\_1 & 100.00 & 15.36 & 37.92 & 2.05 & 6.07 \\
S8/180651 & G06\_2 & 100.00 & 0.87 & 4.45 & 0.77 & 3.83 \\
S8/180651 & G06\_3 & 100.00 & 3.04 & 7.80 & 2.23 & 5.96 \\
S8/180651 & G06\_4 & 99.90 & 5.01 & 13.93 & 12.31 & 35.29 \\
S9/123403 & G07\_1 & 100.00 & 5.00 & 9.78 & 4.49 & 18.84 \\
S9/123403 & G07\_2 & 100.00 & 0.44 & 2.22 & 0.32 & 1.99 \\
S9/123403 & G07\_3 & 100.00 & 1.05 & 3.60 & 2.45 & 8.26 \\
S9/123403 & G07\_4 & 100.00 & 2.80 & 5.69 & 1.90 & 3.75 \\
S9/123725 & G09\_1 & 100.00 & 11.53 & 26.72 & 3.78 & 14.00 \\
S9/123725 & G09\_2 & 100.00 & 6.25 & 16.46 & 4.82 & 17.28 \\
S9/123725 & G09\_3 & 100.00 & 5.51 & 13.62 & 6.89 & 36.18 \\
S9/123725 & G09\_4 & 100.00 & 1.92 & 3.65 & 1.32 & 3.58 \\
S9/124057 & G11\_1 & 100.00 & 1.95 & 6.34 & 2.31 & 4.28 \\
S9/124057 & G11\_2 & 100.00 & 1.11 & 3.09 & 0.91 & 3.17 \\
S9/124057 & G11\_3 & 100.00 & 0.81 & 7.19 & 0.40 & 2.49 \\
S9/124057 & G11\_4 & 100.00 & 12.67 & 28.84 & 4.81 & 8.00 \\
S9/125315 & G21\_1 & 100.00 & 0.83 & 4.39 & 0.64 & 4.58 \\
S9/125315 & G21\_2 & 100.00 & 5.66 & 13.19 & 5.88 & 18.69 \\
S9/125315 & G21\_3 & 100.00 & 5.53 & 12.10 & 1.77 & 4.42 \\
S9/125315 & G21\_4 & 100.00 & 11.16 & 40.09 & 3.89 & 12.78 \\
\bottomrule
\end{longtable}

\begin{longtable}{lp{100mm}}
\caption{Official HO-Cap object identities for the evaluated groups; suffix order is \texttt{\_1} through \texttt{\_4}.}\label{tab:hocap-object-catalog}\\
\toprule Group & Official object descriptions in suffix order \\
\midrule\endfirsthead
\toprule Group & Official object descriptions in suffix order \\
\midrule\endhead
G02 & whole milk carton; cooked ham; chocolate drink powder; sauce bottle \\
G04 & flapjack mix; chocolate fudge; herring fillets can; crackers box \\
G05 & vegetable oil spread; body lotion; peanut chocolate box; sponge \\
G06 & DVD; pumpkin creamer; yellow mustard; dish brush \\
G07 & game controller; hot cocoa mix; dandruff shampoo; toy axe \\
G09 & candy box; toy car; toothpaste box; chocolate syrup bottle \\
G10 & soup mix; gel toothpaste; chocolate biscuit sticks; body wash \\
G11 & coconut water; baby powder; baking soda; chocolate bar \\
G15 & men's body wash; dandelion tea; cooking spray; joy controller \\
G16 & toilet cleaner; laundry detergent; small coconut water; fabric softener sheets \\
G18 & green tea latte mix; projector remote; right shoe; left shoe \\
G19 & electric screwdriver; blue spatula; deodorant; ping pong paddle \\
G20 & cappuccino mix; mustard bottle; toilet cleaner; dog toy bone \\
G21 & moisturizing lotion; playing cards; pink spatula; blue brush \\
G22 & chocolate bar; mayomust sauce; soup mix; gray spatula \\
\bottomrule
\end{longtable}



\end{document}
